\documentclass[11pt,a4paper]{article}

% Essential Packages
\usepackage[utf8]{inputenc}
\usepackage[T1]{fontenc}
\usepackage[english]{babel}
\usepackage{graphicx}
\usepackage{float}
\usepackage{geometry}
\usepackage{amsmath}
\usepackage{amssymb}
\usepackage{rotating}
\usepackage{booktabs}
\usepackage{multirow}
\usepackage{microtype}
\usepackage{siunitx}
\usepackage{adjustbox}
\usepackage[authoryear,round]{natbib}
\usepackage{parskip}
\usepackage{setspace}
\usepackage{tabularx}
\usepackage{lineno}
\onehalfspacing

\IfFileExists{newtxtext.sty}{\usepackage{newtxtext, newtxmath}}{\usepackage{mathptmx}}
\usepackage[font=small, labelfont=bf, margin=1cm]{caption}
\usepackage{authblk}
\usepackage{xurl}
\usepackage{xr}% xr, not xr-hyper: xr-hyper with cleveref printed the file name after every cross-document reference
\usepackage[colorlinks=true, allcolors=blue]{hyperref}
\usepackage{cleveref}
\externaldocument[S-]{supp}
\usepackage{tikz}
\emergencystretch=2em
\usepackage{quantikz}
\usetikzlibrary{shapes.geometric, arrows, positioning, fit, backgrounds}

\usepackage{xcolor}
\newcommand{\tbd}[1]{\textcolor{red}{[#1]}}% pending value: remove before submission

% TODO toggle: \showtodotrue shows the notes, \showtodofalse hides them (submission version)
\newif\ifshowtodo
\showtodofalse
\makeatletter
\newcommand{\todo}{\@ifnextchar\bgroup\@todowitharg\@todonoarg}
\newcommand{\@todowitharg}[1]{\ifshowtodo{\textcolor{red}{\textbf{[TODO: #1]}}}\fi}
\newcommand{\@todonoarg}{\ifshowtodo{\textcolor{red}{\textbf{[TODO]}}}\fi}
\makeatother

% Fail compilation if a figure is missing (prevent silent placeholder boxes)
\graphicspath{{figs/}{./}}
\newcommand{\safefig}[2]{\IfFileExists{#1}{\includegraphics[width=#2]{#1}}{\IfFileExists{figs/#1}{\includegraphics[width=#2]{figs/#1}}{\PackageError{safefig}{figure file \detokenize{#1} not found. Run scripts/make_paper_figs.py first.}{}}}}

% Page Geometry
\geometry{top=2.5cm, bottom=2.5cm, left=2.5cm, right=2.5cm}

\title{\textbf{Limits of Resolution in Small Machine-Learning Benchmarks: A Noise Budget with a Quantum Case Study}}

\author[1,*]{Enrique Real}
\affil[1]{Independent Researcher, Alicante, Spain. ORCID: \href{https://orcid.org/0009-0006-3862-9797}{0009-0006-3862-9797}}
\affil[*]{\textit{Corresponding author: enriquerb32@gmail.com}}
\renewcommand\Affilfont{\itshape\small}

\date{November 2026}

\begin{document}
\maketitle
\linenumbers


\begin{abstract}
Quantum machine learning (QML) benchmarks on small tabular datasets are often read as architecture rankings from AUC differences of a few hundredths. We ask what such a benchmark resolves, on Western Mediterranean monthly precipitation, and report three findings. \emph{Resolution:} with 529 test station-months the standard error of one AUC (about 0.024) is 3--21 times the seed standard deviation of the classical models, the minimum detectable difference between 15-seed means is 0.022--0.032 AUC and falls only as $N^{-1/2}$, and near this limit DeLong and calendar-year cluster procedures disagree. A second four-station panel, a re-test sharing the region and test years, matched the predicted resolution, while the model ranking did not replicate. \emph{Quantum controls:} two hybrid models scored 0.020--0.024 AUC below Random Forest and were not distinguished from classical controls at the same four-angle interface; Fourier surrogates matched the trained circuit, whose readouts are nearly degree-1 tensor polynomials of its angles (also with random weights), and the same controls detected a circuit-generated signal planted at the interface. \emph{References:} no model resolvably improved on a month$\times$station climatology, not even the ECMWF SEAS5 forecast started on the first day of the target month ($+0.043$ AUC, not resolved by the cluster procedures); none discriminated once seasonal and station base rates were removed, and none added information to SEAS5. We propose a reporting protocol: a stated estimand, a minimum detectable difference computed in advance, separate seed and sampling error, base-rate and operational references, and classical controls at the same interface.

\bigskip

\noindent \textbf{Keywords:} Quantum Machine Learning, Benchmarking, Statistical Power, Minimum Detectable Effect, Seed Variance, Precipitation Forecasting.
\end{abstract}

\newpage

% ============================================================
\section{Introduction}
% ============================================================

Claims of quantum utility on classical tabular data frequently rest on a single training run per model, capacity-asymmetric classical baselines, and uncertainty estimates that cover only one source of error \citep{bowles2024}. The machine-learning literature has long pointed out that a single run confounds initialization luck with architectural merit \citep{henderson2018,bouthillier2021}, and that tests computed on one fixed test set ignore the variance induced by the training procedure \citep{dietterich1998,nadeau2003}. Variational circuits aggravate the issue: in the noisy intermediate-scale quantum (NISQ) era \citep{preskill2018}, parameterized quantum circuits are sensitive to initialization even at a few qubits, where barren plateaus \citep{mcclean2018} are not expected and local minima and the periodicity of the encoding matter more, and they are usually combined with a classical backbone that compresses the input into a few qubits \citep{huang2021,schuld2022}.

Throughout, ``noise'' means statistical variability (test-set sampling error and optimization-seed variance), not quantum hardware noise. This paper asks a design question that comes before any ranking: how large a difference in AUC can a benchmark with a few hundred test cases resolve, and how much does the verdict on a small difference depend on the inferential procedure? The answer is largely a property of the evaluation design rather than of the application. The standard error of one AUC is set by the size and prevalence of the test set (about 0.024 for the 529 test cases studied here); it is several times the variation between training seeds of the classical models, and it falls only as $N^{-1/2}$, so repeating training with more seeds does not remove it. For differences of a few hundredths, the paired standard error, the dependence between test cases and the choice of procedure decide the verdict. Variational hybrids add a second complication: their seed variance is large enough that seed and sampling error must both be reported.

We work through these questions on a deliberately modest but realistic case: a multi-decadal, multi-station Western Mediterranean precipitation panel, an area where heavy precipitation is linked to cut-off lows aloft (DANAs; \citealp{romero1999,nieto2005}) and where tree ensembles are strong baselines on engineered tabular features \citep{grinsztajn2022}. The task is a 1-month-ahead probabilistic forecast of precipitation exceeding a pooled regional 60th percentile from antecedent 12-month features. The panel has the features that make small benchmarks hard to read: a few hundred correlated test cases, strong base-rate structure across sites and seasons, and incomplete records. It is a testbed for evaluation methodology, not an operational early-warning system.

Applied climate- and weather-QML studies have addressed separate parts of the problem: opportunities for climate modeling \citep{schwabe2025op}, parameter-count and initialization-matched comparisons for cloud cover \citep{schwabe2025,pastori2026}, wind-speed forecasting \citep{dasilva2026}, quantum kernels for climate emulation \citep{bazgir2024}, synthetic time-series benchmarks \citep{fellner2026}, variational data assimilation \citep{kumar2026}, and parameter-matched neural comparisons \citep{marchisio2026}. In the time-series and weather QML studies cited above, quantum models are reported as better than, or as having the potential to outperform, classical baselines on the basis of point estimates: \citet{marchisio2026} report a 15.4\,\% lower MSE than a parameter-matched classical leaky-integrate-and-fire neuron on a 2{,}000-sample test set, with one value per model and no interval or test, and \citet{dasilva2026} use 10 runs per configuration and show their spread in violin plots, but compare the models through single error values with no interval or test, although their test sets span only 14 and 5 days. We do not claim that these effects are spurious; the point is that the evidence reported cannot separate them from sampling and seed noise.

A second line of work asks what a variational circuit can represent. A circuit that encodes its inputs as rotation angles computes a truncated Fourier series of them, with frequencies fixed by the encoding gates \citep{schuld2021fourier}; such re-uploading models admit classical surrogates \citep{schreiber2023}, random Fourier features approximate many of them efficiently \citep{landman2023}, and necessary and sufficient conditions for this dequantization are known for regression \citep{sweke2025}. More broadly, the structures that make variational models trainable often make them classically simulable \citep{cerezo2025}, quantum kernels concentrate exponentially with the number of qubits \citep{thanasilp2024}, and models that are both trainable and not dequantizable exist but have to be constructed for the purpose \citep{gilfuster2025}. A classical model with the spectrum of the circuit is therefore the natural control for any claimed quantum contribution. Applied benchmarks rarely include it; the case study below does, at the same interface and on the same rows, and checks with a planted signal that the controls can detect a circuit contribution when one exists.

\textbf{What the paper offers to other benchmarks.} The statistical tools are standard (DeLong, cluster bootstrap, minimum detectable difference), and the low power and clustered errors of benchmark comparisons have been documented in NLP \citep{card2020}, neuroimaging \citep{varoquaux2018} and language-model evaluation \citep{miller2024}; the contribution is to combine them around two explicit estimands and to quantify what they imply for benchmark design. Three results are not specific to this panel. First, a resolution rule, checked empirically on the panel: the minimum detectable AUC difference (MDE) of two models compared on the same $N$ test cases falls as $N^{-1/2}$; for 15-seed means of Random Forest against a hybrid it is 0.032 at $N=529$ and, by this scaling, about 0.016 at $N\approx2{,}100$ (\Cref{sec:design}). Second, near this limit verdicts depend on the procedure; known-truth simulations show under which dependence structures each procedure is miscalibrated (\Cref{sec:simulation}), and a decomposition of the cluster-robust variance traces the discrepancy observed on the panel to the co-movement of stations within a year (\Cref{sec:budget}). Third, a reporting protocol whose items need only the test predictions of the models (\Cref{tab:protocol}).

\textbf{What the paper does not claim.} It does not rank the architectures. In particular, it does not show that Random Forest, or tree ensembles in general, beat quantum or neural models on tabular data: no learned model improves on a month$\times$station climatology, and within a station the climatology has the highest AUC (\Cref{sec:not_show}). For the same reason the panel is a poor testbed for choosing an architecture; it is used here for the statistical question, for which small differences between models are the case of interest (\Cref{sec:testbed}).

\textbf{Contributions.}
\begin{enumerate}
    \item \textbf{A noise budget with two explicit estimands}, the difference of single-run AUCs and of seed-averaged AUCs, with paired DeLong, cluster-robust, year-block bootstrap and two-way (years and seeds) bootstrap inference, for three classical and two hybrid quantum models on 529 test station-months, and a diagnosis of the disagreement between procedures (\Cref{sec:framework,sec:budget}).
    \item \textbf{Known-truth simulations} of two correlated scorers with shared and model-specific year effects, which check the type-I error, coverage and power of the procedures at this sample size (\Cref{sec:simulation}).
    \item \textbf{A resolution analysis for benchmark design}, combining a closed-form approximation with an empirical curve of the minimum detectable difference against test size (\Cref{sec:design}).
    \item \textbf{A QML case study and a reporting protocol}: two hybrid quantum models against eleven controls at the same four-dimensional interface (linear, separable and tensor-product Fourier surrogates, frozen circuits and two classical transfer heads), a positive control with a planted circuit signal, a finite-shot simulation of inference, a re-test on a second four-station panel, and climatological and operational (ECMWF SEAS5) references on discrimination, recalibrated proper scoring, within-station discrimination and a label that removes the base rates (\Cref{sec:applying,sec:reference,sec:protocol}).
\end{enumerate}

% ============================================================
\section{Data, Task and Models}
% ============================================================

\subsection{Dataset, Target and Split}
\label{sec:dataset_merge}

The study domain is the Spanish Mediterranean coast (Catalonia, the Valencian Community and the Region of Murcia), a transition zone between Mediterranean (\textit{Csa}) and semi-arid (\textit{BSk}/\textit{BSh}) K\"{o}ppen regimes subject to severe convective episodes \citep{romero1999,nieto2005}. Daily records of four AEMET stations, Alacant/Alicante, Barcelona Aeropuerto, Val\`{e}ncia and San Javier Aeropuerto (\Cref{tab:station_coverage}), were obtained from \textit{datosclima.es} under the AEMET Open Data framework (downloaded on 1 June 2026) and aggregated to calendar months per station. The \emph{label} is computed from the precipitation series alone: the monthly total is defined when at most 3 days of the month are missing (trace precipitation counts as zero) and is otherwise left undefined, so no month is labelled from an incomplete record or from days dropped because another variable was missing. The \emph{predictors} are the precipitation totals of the 12 preceding months (rescaled to the month length when at least 10 days are observed) and monthly means over the observed days of pressure at station level and its variability, temperature and diurnal range, wind speed and maximum gust, sunshine duration, and FAO-56 Penman--Monteith reference evapotranspiration \citep{allen1998} with humidity estimated from the minimum temperature, radiation from sunshine duration and wind from the observed speed ($u_2=0.748\,u_{10}$). A covariate needs at least 10 observed days in a month; gaps of up to two months are interpolated within a station, and longer gaps leave the month without predictors. Every calendar month is kept, so the windows are contiguous in time, and a window is used only if its target month has a label, its last antecedent month has all predictors and its 12 precipitation lags are defined. Inhomogeneities from relocations or instrument changes remain an accepted limitation \citep{alexandersson1986,peterson1998}.

Each window uses months $i,\dots,i+11$ (35 features: 12 precipitation lags, 17 engineered physical features, 2 cyclic month terms and 4 station indicators) to predict whether precipitation in month $i+12$ exceeds the pooled regional 60th percentile of the training labels, $P_{60}=\qty{29.9}{\milli\meter}$. The split is chronological: training 1938--1996 (2{,}332 windows), validation 1997--2007 (516) and test 2008--2026 (529; 40.6\,\% positive). The precipitation series are almost complete (0--5\,\% missing days), but the pressure records of Valencia and San Javier are not (19 and 25\,\% of days missing over their spans), and after 2012 and 2011 these two stations have no complete set of predictors; 83.4\,\% of the test windows therefore come from Alicante and Barcelona. Test base rates differ between stations (27--55\,\%), so part of the pooled AUC reflects telling stations apart (\Cref{sec:reference}). All scaling and anomaly statistics use pre-1997 data only. The test period falls in a time of accelerated warming, so the evaluation takes place under a covariate shift whose effect cannot be isolated with one chronological split.

\textbf{Revision of the data pipeline.} An earlier version of this analysis built the monthly panel from an inner join of the five daily variable files, so that a day without pressure or wind also removed its precipitation; 23\,\% of the test months were then labelled from incomplete totals. It also set missing precipitation and sunshine to zero, used a constant wind in the evapotranspiration, and joined non-contiguous months into windows. The results reported here come from the corrected pipeline, rerun from the raw files with tuned tree baselines in one pinned software environment; \Cref{S-app:v1} lists the changes and compares the headline numbers of the two versions. The test set changed from 528 to 529 windows with the same station composition, so the missing stations after 2012 are a property of the covariate records, not of the join.

\begin{table}[htbp]
\centering
\caption{AEMET stations (climatological identifier, altitude), test windows and test-set class prevalence for the regional exceedance threshold ($P_{60}=29.9$\,mm). Coordinates: Alicante \ang{38;22;21}\,N \ang{0;29;39}\,W; Barcelona \ang{41;17;30}\,N \ang{2;3;38}\,E; San Javier \ang{37;47;20}\,N \ang{0;48;12}\,W; Valencia \ang{39;28;50}\,N \ang{0;21;59}\,W. Last test year: the last year with a complete set of predictors. Missing days of the precipitation and pressure files over the span of each station in \Cref{S-tab:etl}.}
\label{tab:station_coverage}
{\small
\begin{tabular}{llccc}
\toprule
\textbf{Station} & \textbf{AEMET id (alt.)} & \textbf{Test windows} & \textbf{Last test year} & \textbf{Test positive rate} \\
\midrule
Alicante    & 8025 (81\,m) & 220 & 2026 & 27.3\% (60 / 220) \\
Barcelona   & 0076 (1\,m)  & 221 & 2026 & 51.1\% (113 / 221) \\
San Javier  & 7031 (4\,m)  & 37  & 2011 & 37.8\% (14 / 37) \\
Valencia    & 8416 (13\,m) & 51  & 2012 & 54.9\% (28 / 51) \\
\midrule
\textbf{Total} & & \textbf{529} & & \textbf{40.6\% (215 / 529)} \\
\bottomrule
\end{tabular}%
}
\end{table}

\subsection{Models and Fitting Sets}
\label{sec:architectures}

\textbf{Models in the main analyses.} Two references, the ECMWF SEAS5 seasonal forecast (\Cref{sec:seas5}) and a month$\times$station climatology fitted on the training partition, are compared with three classical models, Random Forest \citep{breiman2001}, XGBoost and a standalone bidirectional LSTM (10{,}103 parameters), and with two hybrid quantum--classical models simulated without noise in PennyLane (state-vector simulator \texttt{default.qubit}, exact expectation values, gradients by backpropagation, double precision). The two tree ensembles were tuned by an automated search on the validation AUC with declared search spaces (Optuna, TPE sampler, 50 trials each): Random Forest selected 350 trees of depth 11 (minimum split 13, minimum leaf 10, $\sqrt{p}$ features), and XGBoost, which stops early on the validation AUC, a learning rate of 0.036 with trees of depth 3; the search spaces, the selected values and the untuned configurations are in \Cref{S-app:hyperparams}. \texttt{LSTM+QNN} compresses a 51-dimensional vector (the 16-dimensional LSTM state concatenated with the 35 features) to 4 rotation angles $q\in[0,\pi]^4$ that feed a 4-qubit circuit of two layers, each re-uploading the angles as $R_Z(q_j)R_Y(q_j)$ ($R_Y$ acts first on $|0\rangle$, so $R_Z$ is not a global phase) before one \texttt{StronglyEntanglingLayers} block of range $r=1$ (a ring of nearest-neighbour CNOTs; the template is called once per layer, so both layers have $r=1$); the readout is $\langle Z_j\rangle$ on every qubit, followed by LayerNorm and a linear output layer. Because LayerNorm standardizes the four readouts of each row, the model is invariant to a common contraction of the $\langle Z_j\rangle$, such as global depolarizing noise would cause, but not to shot noise or to qubit-dependent noise. Each angle enters four rotation gates, so the model computes a truncated Fourier series of the angles with per-coordinate frequencies $\{0,\pm1,\dots,\pm4\}$, with products across coordinates generated by the entangling gates \citep{schuld2021fourier}; it has 2{,}613 parameters, 24 of them quantum. \texttt{LSTM+QTL} is a quantum head on a pre-trained standalone LSTM \citep{mari2020}: in each seed the LSTM of that seed is frozen in evaluation mode, and a trainable adapter ($67\to8\to4$, angles in $[0,\pi]$) feeds a 2-layer circuit with a single $R_Y$ encoding (per-coordinate frequencies $\{0,\pm1\}$; one \texttt{StronglyEntanglingLayers} call with two layers, so PennyLane's default ranges $r=1$ and $r=2$). A variant of \texttt{LSTM+QNN} trained with a hand-set penalty (\texttt{LSTM+PI-QNN}) is reported in \Cref{S-app:hybrids}.

\textbf{Controls at the interface.} Eleven classical or frozen-circuit controls isolate the circuit. The Linear-4D control has the architecture of \texttt{LSTM+QNN} with the circuit replaced by a linear $4\to4$ map (2{,}609 parameters against 2{,}613). Four separable Fourier controls apply the linear map to $[\sin kq,\cos kq]_{k\le d}$ for $d=1,\dots,4$ (\texttt{Fourier-1} to \texttt{Fourier-4}; 36--132 parameters in the map), the last one covering the circuit's full per-coordinate spectrum without its cross-coordinate products. Two tensor-product surrogates, \texttt{Full Fourier-1} and \texttt{Full Fourier-2}, use the bases $\prod_{j=1}^{4}\{1,\cos k q_j,\sin k q_j\}_{k\le d}$ for $d=1$ (81 features) and $d=2$ (625 features); a full degree-4 basis would have $9^4=6{,}561$ features for 2{,}332 training rows. Two controls keep the circuit but freeze its weights, at random values in $(-\pi,\pi)$ or at zero (the encodings and the CNOT rings remain, so the circuit is a fixed nonlinear map of the angles). For \texttt{LSTM+QTL}, \texttt{QTL-classical} replaces its circuit by a linear $4\to4$ map, and \texttt{QTL tensor-81} by a linear map on the 81 functions $\prod_j\{1,\cos q_j,\sin q_j\}$, which span exactly the readouts of a circuit with one $R_Y$ encoding per angle; it is the exact classical counterpart of that head. All controls share the training loop of the model they are compared with (\Cref{S-tab:core_spec}). Comparisons of a hybrid with the standalone LSTM, which has about four times its parameters, or with Random Forest are comparisons of whole pipelines and cannot be attributed to the circuit; the comparisons that isolate the circuit are those with the controls.

\textbf{Fitting sets, seeds and environment.} All models use the same 2{,}332 training rows; the neural and hybrid models and XGBoost use the 516 validation rows for early stopping; all are scored on the same 529 test rows, with 15 seeds each. All results come from one execution in a pinned environment (Python 3.12, PyTorch 2.15.0 CPU, PennyLane 0.45.1, scikit-learn 1.9.1, XGBoost 2.0.3; lock file with hashes in the repository; \Cref{S-tab:software_env}) with a fixed number of threads. The primary metric is the ROC-AUC, complemented by the Brier score \citep{brier1950}.

% ============================================================
\section{Statistical Framework}
\label{sec:framework}
% ============================================================

\subsection{Estimands and Inference}
\label{sec:estimand}

For models $A$ and $B$ with test probabilities $p_{A,s}$, $p_{B,s}$ from seeds $s=1,\dots,S$ ($S=15$), the main test statistic of this paper is the difference between seed-averaged test AUCs,
\begin{equation}
\hat\Delta=\frac1S\sum_{s}\widehat{\mathrm{AUC}}(p_{A,s})-\frac1S\sum_{s}\widehat{\mathrm{AUC}}(p_{B,s}),
\end{equation}
and, for single-run comparisons, the difference of two individual AUCs. As in forecast verification, the two models are scored on the same cases, so their errors are dependent and must be compared paired \citep{diebold1995,delsole2014}. Sampling inference targets the corresponding performance difference for the evaluation population, conditional on the observed training data and on the fitted seed-specific models. Variation due to the choice of training set, which can dominate comparisons of learning algorithms \citep{nadeau2003,bouthillier2021} and is what cross-validation estimates \citep{bates2024}, is outside this estimand: the chronological split is fixed, as it is in operational use. The DeLong variance conditions on the observed test years and treats station-months as exchangeable; the year-block bootstrap and the simulations of \Cref{sec:simulation} instead target the difference for new years. A bootstrap that also resamples seeds targets the seed-averaged difference over the distribution of optimization seeds. The mean of AUCs is linear in the DeLong structural components, so the paired sampling variance of $\hat\Delta$ is obtained by averaging the components over seeds; this treats the seeds as fixed.

\subsection{Sampling Uncertainty}

Test-set sampling error is estimated in four ways. (i) The paired DeLong estimator for correlated ROC curves \citep{delong1988,hanley1982} treats station-months as independent. (ii) A calendar-year cluster bootstrap (``year-block bootstrap'') resamples the 19 test years (2008--2026) with all stations of a year together ($B=2{,}000$; \citealp{field2007}); intervals are percentile intervals. (iii) A two-way bootstrap resamples years and, within each model, seeds, independently between models. (iv) An analytic cluster-robust estimator \citep{obuchowski1997} sums the influence values of $\hat\Delta$ (the seed-averaged DeLong components, centered and divided by the number of positives or negatives) within each year, $\text{SE}_{\text{cl}}=\{K/(K-1)\sum_u S_u^2\}^{1/2}$ with $S_u$ the sum over the rows of year $u$; DeLong is the same sum without grouping, and the ratio $\text{SE}_{\text{cl}}/\text{SE}_{\text{DeLong}}$ (``cluster/iid'') summarizes the within-year dependence. Splitting $S_u$ into the contributions of the stations separates the cluster variance into within-station and cross-station terms, and permuting the years independently within each station gives its distribution when stations do not co-move within a year. With only 19 clusters, cluster-based variance estimates are biased downward \citep{cameron2008}; the Bell--McCaffrey CR2 correction \citep{bell2002} and a studentized wild cluster bootstrap with Webb weights \citep{webb2023} are reported as checks (\Cref{S-app:fewclusters}). All procedures are reported, and none is assumed a priori to be the correct one for the observational panel.

\subsection{Seed Variability, Minimum Detectable Difference and Scope}

For a model with seed standard deviation $\sigma_{\text{seed}}$ and mean sampling standard error $\text{SE}$ of one AUC we report the seed share $\sigma_{\text{seed}}^2/(\sigma_{\text{seed}}^2+\text{SE}^2)$, a descriptive ratio of two conditional variances, not a decomposition of the total uncertainty of a training procedure. For a paired difference with standard error $\text{SE}_\Delta$, the minimum detectable difference at two-sided $\alpha=0.05$ and power 0.80 is $\text{MDE}=(z_{0.975}+z_{0.80})\,\text{SE}_\Delta\approx2.80\,\text{SE}_\Delta$. It describes the sensitivity of the design to a true difference; it is not a threshold that an observed difference must exceed, and an observed gap between $1.96\,\text{SE}_\Delta$ and $2.80\,\text{SE}_\Delta$ can be significant.
\label{sec:multiplicity}

The primary analyses are the noise budget, the paired comparisons, the simulations and the comparison with the climatology; none was pre-registered. Holm correction is applied to the DeLong $p$-values within two families of three seed-mean comparisons (the classical models, and \texttt{LSTM+QNN} against Random Forest, the Linear-4D control and the LSTM). The remaining analyses, about forty paired comparisons, are secondary, and isolated nominal rejections among them are expected by chance. One rule is used for every comparison, primary or secondary: a difference is called \emph{resolved} when the paired DeLong test, the calendar-year block bootstrap and the cluster-robust interval ($t_{18}$) all exclude zero, and otherwise we state which procedures do.

% ============================================================
\section{Results}
% ============================================================

\subsection{The Noise Budget}
\label{sec:budget}

On the scale of one AUC, test-sampling variability exceeds seed variability for every model (\Cref{tab:budget}). One AUC has a standard error of 0.023--0.024 under DeLong and 0.026--0.027 under the year-block bootstrap, whereas the seed standard deviation is 0.001--0.008 for the classical models (seed share 0.2--11\,\%) and 0.010--0.014 for the hybrids (14--25\,\%). For variational models the seed term is therefore not small, and both sources must be reported.

\begin{table}[htbp]
\centering
\caption{Noise budget for one AUC (15 seeds per model, 529 test rows). iid: DeLong; block: calendar-year block bootstrap ($B=2{,}000$, 19 years). Seed share: $\sigma_{\text{seed}}^2/(\sigma_{\text{seed}}^2+\text{SE}^2)$, with the DeLong or the block standard error.}
\label{tab:budget}
{\small
\adjustbox{max width=\textwidth}{%
\begin{tabular}{lccccc}
\toprule
\textbf{Quantity} & \textbf{Random Forest} & \textbf{XGBoost} & \textbf{Standalone LSTM} & \textbf{LSTM+QNN} & \textbf{LSTM+QTL} \\
\midrule
Mean AUC (15 seeds) & 0.6784 & 0.6684 & 0.6638 & 0.6587 & 0.6548 \\
SE of one AUC, iid / block & 0.023 / 0.026 & 0.024 / 0.027 & 0.024 / 0.027 & 0.024 / 0.026 & 0.024 / 0.027 \\
SD across seeds & 0.0011 & 0.0081 & 0.0051 & 0.0139 & 0.0097 \\
\textbf{Seed share, iid / block} & \textbf{0.2\,\%} / 0.2\,\% & \textbf{10.6\,\%} / 8.2\,\% & \textbf{4.5\,\%} / 3.4\,\% & \textbf{25.2\,\%} / 21.6\,\% & \textbf{14.2\,\%} / 11.3\,\% \\
\bottomrule
\end{tabular}}%
}
\end{table}

\Cref{tab:paired_ci} and \Cref{fig:paired_ci} compare 15-seed means. Pairing halves the standard error, and averaging seeds lowers it further, but the DeLong MDE of the classical comparisons remains 0.022--0.031. No classical difference is resolved, and none survives the Holm correction (smallest $p=0.186$). For RF $-$ LSTM ($+0.015$) the procedures disagree: DeLong gives $p=0.186$, whereas the year-block interval ($[+0.002,+0.028]$), the two-way interval and the cluster-robust estimator (standard error 0.0065 against 0.0110; CR2 $p=0.040$, wild cluster bootstrap $p=0.041$) exclude zero. A single run is less informative still: with one seed per model the paired DeLong MDE is 0.031--0.044, and Random Forest is declared better than XGBoost, the LSTM or \texttt{LSTM+QNN} in 7\,\%, 0\,\% and 9\,\% of the $15\times15$ seed pairings (\Cref{S-app:single_seed}).

\begin{table}[htbp]
\centering
\caption{Paired comparison of 15-seed means of the classical models: $\hat\Delta$ with the paired DeLong standard error on seed-averaged components, two-sided $p$-value, 95\,\% intervals from the calendar-year block bootstrap and from the years$\times$seeds bootstrap ($B=2{,}000$), ratio of the cluster-robust (calendar-year) to the DeLong standard error, and DeLong MDE ($2.80\times$SE). Holm first-step threshold 0.017.}
\label{tab:paired_ci}
{\small
\adjustbox{max width=\textwidth}{%
\begin{tabular}{lccccccc}
\toprule
Difference & $\hat\Delta$ & SE (DeLong) & $z$ ($p$) & Year-block 95\,\% CI & Years$\times$seeds 95\,\% CI & Cluster/iid & MDE \\
\midrule
RF $-$ XGBoost & $+0.0100$ & 0.0079 & $1.27$ (0.205) & $[-0.002,+0.020]$ & $[-0.004,+0.021]$ & 0.75 & 0.022 \\
RF $-$ LSTM & $+0.0146$ & 0.0110 & $1.32$ (0.186) & $[+0.002,+0.028]$ & $[+0.001,+0.029]$ & 0.59 & 0.031 \\
XGBoost $-$ LSTM & $+0.0046$ & 0.0104 & $0.44$ (0.662) & $[-0.010,+0.023]$ & $[-0.011,+0.024]$ & 0.77 & 0.029 \\
\bottomrule
\end{tabular}}%
}
\end{table}

\begin{figure}[tbp]
    \centering
    \safefig{figs/fig_forest.pdf}{0.78\textwidth}
    \caption{Paired differences of 15-seed mean test AUC of \Cref{tab:paired_ci,tab:qnn_paired,tab:controls}, with the 95\,\% intervals of three procedures: DeLong (circles), calendar-year block bootstrap (squares) and years$\times$seeds bootstrap (triangles); 529 test rows, $B=2{,}000$. Intervals are pointwise and not multiplicity-adjusted.}
    \label{fig:paired_ci}
\end{figure}

\textbf{Why the procedures disagree.} The disagreement appears for Random Forest against neural and hybrid models (cluster/iid 0.59--0.78) and not among neural and hybrid models (0.81--1.42). Writing the cluster variance as the sum of within-station and cross-station terms locates it (\Cref{S-app:cancellation}): for RF $-$ LSTM the cross-station term is $-42$\,\% of the within-station one, that is, in years in which Random Forest gains over the LSTM at one station it tends to lose at another, and the yearly contributions of the stations are negatively correlated (mean $-0.15$). Permuting the years independently within each station, which removes this co-movement and keeps everything else, gives cluster/iid ratios with median 0.75 for RF $-$ LSTM and 0.87 for RF $-$ \texttt{LSTM+QNN}; the observed ratios (0.59 and 0.72) lie in the lower tail ($p=0.07$ and 0.12). Removing a constant per station does not change the ratios (0.62 and 0.69), whereas removing the mean of each station$\times$calendar month$\times$class cell restores them to 1.0--1.06: the cancellation comes from station-specific seasonal patterns in which the two model families differ, which offset each other across stations within the same year. The earlier version of the analysis, with labels affected by the inner join, gave ratios of 0.55--0.71 and a significant permutation test (\Cref{S-app:v1}); with corrected labels the effect is weaker and consistent with chance co-movement of four stations over 19 years, and it did not appear on the second panel (\Cref{sec:replication}). The cluster procedures are therefore not wrong on this panel, but their advantage over DeLong rests on a structure of the test years that need not recur. The usual small-sample corrections do not change the cluster verdicts (CR2 equals CR1 to three decimals; \Cref{S-tab:fewclusters}).

\textbf{Re-execution.} Neural training on CPU is not bit-reproducible across software versions and thread settings. Within the pinned environment, refitting the 15 \texttt{LSTM+QNN} seeds for the finite-shot analysis gave a mean AUC of 0.6599 against 0.6587, a change of 0.001, below the standard deviation of a 15-seed mean ($0.014/\sqrt{15}=0.0036$) and an order of magnitude below the MDE.

\subsection{Simulation with Known Truth}
\label{sec:simulation}

Because DeLong and the cluster procedures disagree on gaps near the MDE, we check both on data whose true difference is known. We simulate two scorers whose population AUC difference is known by construction or approximated on large reference samples. Labels have prevalence 0.41 with year-level effects on the logit, and 528 rows are split among 19 years. The scorers are binormal with noise correlation $\rho=0.9$, which reproduces the paired DeLong standard error of the classical comparisons (0.011), and their design AUCs without year effects are 0.65 and $0.65+\Delta$. Year effects enter as a prevalence shift ($\tau_\pi$), a shift shared by both scorers ($\tau_s$) and, in scenarios C--E, a shift specific to each scorer ($\tau_m$). Scenario B has $\tau_\pi=\tau_s=0.3$, C adds $\tau_m=0.3$, D is C with unequal year sizes, and E has all three equal to 0.6 (details in \Cref{app:noise_budget}). Each scenario has 500 replications.

The simulations (\Cref{S-tab:sim,S-fig:sim}) give three results. First, with year effects shared by the two scorers (scenario B) DeLong is calibrated (type-I error 0.040, coverage 0.960), whereas the block bootstrap rejects the exact null in 7.4\,\% of replications. Second, with scorer-specific year effects the DeLong standard error is too small (ratio of Monte Carlo SD to mean SE 1.07, 1.17 and 1.57 in scenarios C--E; type-I error 0.068, 0.096 and 0.220), whereas the block standard error tracks the SD (1.01--1.08) with a type-I error of 0.074--0.084: neither procedure is uniformly better calibrated. Third, where DeLong is calibrated power follows the MDE rule: in scenario B (MDE 0.031) a true difference of 0.029 is detected in 74\,\% of replications and 0.0145 in 32\,\%. A version calibrated to the year sizes, prevalence and score correlation of the panel gives type-I errors of 0.048--0.058 for DeLong and 0.064--0.088 for the block bootstrap, and only 17--24\,\% power at an observed gap of 0.013 (\Cref{S-app:panel_matched}). None of these mechanisms produces a cluster standard error well below the DeLong one; that pattern needs the cross-station co-movement described in \Cref{sec:budget}.

\subsection{Case Study: Hybrid Quantum Models under the Noise Budget}
\label{sec:applying}

\Cref{tab:case} lists the 15-seed AUCs of all models fitted on the same rows. The climatology has the highest value; Random Forest is the best learned model, and every hybrid lies below it and below the standalone LSTM, with seed standard deviations of 0.010--0.020.

\begin{table}[htbp]
\centering
\caption{15-seed test AUC (same 2{,}332 fitting rows; mean and sample SD; 529 test rows). Last column: difference to \texttt{LSTM+QNN}, descriptive. Controls in the bottom block, sorted by mean; in brackets the size of the map that replaces the circuit. The extractor of \texttt{LSTM+QTL}, \texttt{QTL-classical} and \texttt{QTL tensor-81} is the standalone LSTM of the same seed.}
\label{tab:case}
{\small
\begin{tabular}{lccc}
\toprule
\textbf{Model} & \textbf{Mean AUC} & \textbf{SD} & \textbf{Gap to LSTM+QNN} \\
\midrule
Month$\times$station climatology & 0.6851 & -- & $+0.0264$ \\
\midrule
Random Forest (tuned) & 0.6784 & 0.0011 & $+0.0197$ \\
XGBoost (tuned, early stopping) & 0.6684 & 0.0081 & $+0.0097$ \\
Standalone LSTM & 0.6638 & 0.0051 & $+0.0051$ \\
LSTM+QNN (24 quantum parameters) & 0.6587 & 0.0139 & $+0.0000$ \\
LSTM+QTL & 0.6548 & 0.0097 & $-0.0039$ \\
LSTM+PI-QNN (penalized variant) & 0.6539 & 0.0201 & $-0.0048$ \\
\midrule
Full Fourier-1 (81 features) & 0.6671 & 0.0054 & $+0.0084$ \\
Fourier-1 (36) & 0.6638 & 0.0055 & $+0.0051$ \\
Fourier-2 (68) & 0.6628 & 0.0037 & $+0.0041$ \\
Fourier-3 (100) & 0.6627 & 0.0058 & $+0.0040$ \\
Fourier-4 (132) & 0.6619 & 0.0069 & $+0.0032$ \\
QTL tensor-81 (linear map on 81 features) & 0.6615 & 0.0062 & $+0.0028$ \\
Full Fourier-2 (625 features) & 0.6613 & 0.0063 & $+0.0026$ \\
QTL-classical (linear head) & 0.6575 & 0.0055 & $-0.0012$ \\
Linear-4D control (20) & 0.6553 & 0.0040 & $-0.0034$ \\
Frozen random circuit & 0.6466 & 0.0201 & $-0.0121$ \\
Frozen zero-angle circuit & 0.6458 & 0.0167 & $-0.0129$ \\
\bottomrule
\end{tabular}%
}
\end{table}

\begin{table}[htbp]
\centering
\caption{Paired comparisons of the hybrids (15-seed means; procedures and notation as in \Cref{tab:paired_ci}). The first three rows form the Holm family (first-step threshold 0.017); the others are secondary.}
\label{tab:qnn_paired}
{\small
\adjustbox{max width=\textwidth}{%
\begin{tabular}{@{}lccccccc@{}}
\toprule
Difference & $\hat\Delta$ & SE (DeLong) & $z$ ($p$) & Year-block 95\,\% CI & Years$\times$seeds 95\,\% CI & Cluster/iid & MDE \\
\midrule
RF $-$ LSTM+QNN & $+0.0197$ & 0.0115 & $1.71$ (0.088) & $[+0.005,+0.038]$ & $[+0.002,+0.041]$ & 0.72 & 0.032 \\
Linear-4D $-$ LSTM+QNN & $-0.0034$ & 0.0059 & $-0.58$ (0.565) & $[-0.016,+0.008]$ & $[-0.020,+0.011]$ & 1.09 & 0.016 \\
LSTM $-$ LSTM+QNN & $+0.0051$ & 0.0057 & $0.90$ (0.368) & $[-0.004,+0.015]$ & $[-0.009,+0.019]$ & 0.88 & 0.016 \\
RF $-$ LSTM+QTL & $+0.0235$ & 0.0100 & $2.35$ (0.019) & $[+0.009,+0.040]$ & $[+0.006,+0.042]$ & 0.78 & 0.028 \\
LSTM+QTL $-$ its extractor & $-0.0090$ & 0.0049 & $-1.85$ (0.065) & $[-0.018,+0.001]$ & $[-0.021,+0.004]$ & 1.04 & 0.014 \\
LSTM+QTL $-$ QTL-classical & $-0.0027$ & 0.0041 & $-0.67$ (0.505) & $[-0.010,+0.006]$ & $[-0.014,+0.009]$ & 1.00 & 0.011 \\
LSTM+QTL $-$ QTL tensor-81 & $-0.0066$ & 0.0042 & $-1.60$ (0.111) & $[-0.014,+0.001]$ & $[-0.019,+0.006]$ & 0.93 & 0.012 \\
\bottomrule
\end{tabular}}}
\end{table}

\textbf{Against classical pipelines.} Random Forest exceeds \texttt{LSTM+QNN} by 0.020 (\Cref{tab:qnn_paired}). As for RF $-$ LSTM, the procedures disagree: DeLong does not reject the null ($p=0.088$; Holm-adjusted 0.26), whereas the year-block, two-way and cluster-robust intervals exclude zero (cluster-robust $t_{18}$ $[+0.002,+0.037]$, CR2 $p=0.029$, wild cluster bootstrap $p=0.022$). Random Forest exceeds \texttt{LSTM+QTL} by 0.024, and this secondary comparison is resolved by every procedure ($p=0.019$; year-block $[+0.009,+0.040]$; cluster-robust $[+0.007,+0.040]$). These are comparisons of whole pipelines: the standalone LSTM, with the same inputs and no bottleneck, lies between Random Forest and the hybrids and is not distinguished from \texttt{LSTM+QNN} ($+0.005$, MDE 0.016).

\textbf{Against controls at the interface.} In the capacity-matched comparison the hybrid is not distinguished from the circuit-free Linear-4D control ($-0.003$ for the control; MDE 0.016), and \texttt{LSTM+QTL} is not distinguished from the same model with a linear head ($-0.003$, MDE 0.011) or with the 81-function tensor head that spans its circuit ($-0.007$, $p=0.11$, MDE 0.012); it scores 0.009 below its own frozen extractor ($p=0.065$).

\begin{table}[htbp]
\centering
\caption{Paired comparisons of the trained circuit of \texttt{LSTM+QNN} with Fourier surrogates and frozen circuits (15-seed means; procedures as in \Cref{tab:paired_ci}). Positive $\hat\Delta$: the control scores above the circuit. Secondary analysis: $p$-values unadjusted.}
\label{tab:controls}
{\small
\adjustbox{max width=\textwidth}{%
\begin{tabular}{@{}lccccccc@{}}
\toprule
Difference & $\hat\Delta$ & SE (DeLong) & $z$ ($p$) & Year-block 95\,\% CI & Years$\times$seeds 95\,\% CI & Cluster/iid & MDE \\
\midrule
Fourier-1 $-$ LSTM+QNN & $+0.0051$ & 0.0045 & $1.12$ (0.262) & $[-0.004,+0.015]$ & $[-0.009,+0.018]$ & 1.15 & 0.013 \\
Fourier-2 $-$ LSTM+QNN & $+0.0041$ & 0.0046 & $0.88$ (0.381) & $[-0.005,+0.015]$ & $[-0.009,+0.018]$ & 1.12 & 0.013 \\
Fourier-3 $-$ LSTM+QNN & $+0.0040$ & 0.0048 & $0.84$ (0.400) & $[-0.006,+0.016]$ & $[-0.009,+0.020]$ & 1.20 & 0.013 \\
Fourier-4 $-$ LSTM+QNN & $+0.0032$ & 0.0050 & $0.63$ (0.526) & $[-0.009,+0.017]$ & $[-0.013,+0.022]$ & 1.42 & 0.014 \\
Full Fourier-1 $-$ LSTM+QNN & $+0.0084$ & 0.0045 & $1.85$ (0.064) & $[-0.002,+0.020]$ & $[-0.005,+0.023]$ & 1.24 & 0.013 \\
Full Fourier-2 $-$ LSTM+QNN & $+0.0026$ & 0.0048 & $0.53$ (0.597) & $[-0.008,+0.014]$ & $[-0.012,+0.019]$ & 1.20 & 0.014 \\
Frozen random $-$ LSTM+QNN & $-0.0121$ & 0.0038 & $-3.15$ (0.002) & $[-0.020,-0.003]$ & $[-0.028,+0.002]$ & 1.15 & 0.011 \\
Frozen zero-angle $-$ LSTM+QNN & $-0.0129$ & 0.0041 & $-3.18$ (0.001) & $[-0.022,-0.005]$ & $[-0.030,+0.004]$ & 1.08 & 0.011 \\
Full Fourier-2 $-$ Fourier-2 & $-0.0015$ & 0.0030 & $-0.50$ (0.615) & $[-0.006,+0.003]$ & $[-0.009,+0.006]$ & 0.76 & 0.008 \\
\bottomrule
\end{tabular}}}
\end{table}

\textbf{Is it the training, and is it quantum?} The controls of \Cref{tab:controls} address two questions. Freezing the circuit weights, at random values or at zero, lowers the AUC by 0.012 and 0.013, a difference resolved by DeLong, the year-block bootstrap and the cluster-robust interval; the frozen models also vary more between seeds (SD 0.017--0.020), and the two-way interval, which includes that variation, reaches zero ($[-0.028,+0.002]$ and $[-0.030,+0.004]$). Training the circuit therefore helps, conditional on the fitted seeds. Whether it helps because the map is quantum is answered by the surrogates: the four separable Fourier controls score 0.003--0.005 above the trained circuit, \texttt{Full Fourier-2} 0.003 above it and \texttt{Full Fourier-1}, the degree-1 tensor basis, 0.008 above it ($p=0.064$; every interval includes zero), all with MDEs of 0.013--0.014. The 90\,\% cluster-robust intervals bound the difference between the circuit and each surrogate within $\pm0.013$--$0.018$ AUC (\Cref{S-tab:fewclusters}); these are descriptive bounds, not equivalence tests, since no margin was fixed in advance. On this panel, then, nothing that the trained circuit adds requires a quantum map: a classical periodic feature map of the same four angles does as well, within a resolution of about 0.015 AUC. This is the empirical counterpart of the classical-surrogate argument \citep{schreiber2023,landman2023}, with a different control: the surrogates here are separate model classes fitted end to end like the hybrid, whereas those of \citet{schreiber2023} are derived from a given circuit.

\textbf{What the circuit computes.} A least-squares fit of the degree-1 tensor basis to the exact readouts of the 15 trained circuits on the test angles gives $R^2$ of 0.95--1.00 (mean 0.984), against 0.85 for separable bases; on the same angles, circuits with random weights give 0.97 (\Cref{S-app:spectrum}). For the median seed, the central 98\,\% of the visited angles of each qubit lie between about 0.4--1.0 and 2.3--2.6\,rad, and in some seeds up to 14\,\% lie within $\pi/18$ of the ends of $[0,\pi]$. The near-linearity in each coordinate is therefore a property of this encoding on this angle range, not something the training finds; and 68\,\% of the readout variance of the trained circuits is at frequency 1 per coordinate and 1\,\% at frequency 4.

\textbf{Can the controls detect a circuit contribution?} With no learned model above the climatology, ``the surrogates match the circuit'' could hold simply because there is nothing for the circuit to learn. We therefore planted one (\Cref{S-app:teacher}). On the real inputs and splits, teacher angles were computed from a fixed random projection, a circuit with the architecture of \texttt{LSTM+QNN} and fixed random weights produced a score, and labels were drawn from the score with noise set for an oracle AUC of 0.85; among 20 random teachers the one with the largest non-separable part was kept (89\,\% of its variance outside a separable degree-4 basis). Given the teacher angles, the trained circuit reaches AUC 0.706 (5 seeds), clearly above the Linear-4D map (0.613), \texttt{Fourier-2} (0.655) and \texttt{Fourier-4} (0.677), and the tensor surrogates reach 0.804 (\texttt{Full Fourier-1}) and 0.810 (\texttt{Full Fourier-2}), near the oracle. The controls thus separate a circuit contribution from a linear or separable map when one exists, and the tensor surrogate reproduces it, as the dequantization results predict for a circuit of this size. When the whole models must also learn the projection from the raw inputs, none of them recovers the planted signal (AUC 0.59--0.62 for every model, including the hybrid); a comparison of complete models on this panel would not detect a circuit contribution of this kind, which is why the interface controls are the relevant ones.

\textbf{Finite shots.} On hardware the noise budget gains a third term: each $\langle Z_j\rangle$ estimated from $n_{\text{shots}}$ measurements adds per-row noise. Sampling the 16-outcome distribution of each of the 15 circuits 50 times per budget (\Cref{S-app:shots}), the 15-seed mean AUC falls by 0.029, 0.012, 0.002 and 0.0001 at 32, 100, 1{,}000 and 10{,}000 shots per evaluation, with a shot standard deviation of the seed-mean AUC of 0.003, 0.002, 0.001 and 0.0004. From about 1{,}000 shots the third term is an order of magnitude below the seed and sampling terms; at 32 shots the loss approaches the MDE. Training on hardware was not simulated: with parameter-shift gradients each row and step needs 81 circuit evaluations.

\subsection{Reference Baselines: Climatology, Proper Scores and Base Rates}
\label{sec:reference}

The comparisons so far are between learned models. A stricter question is whether any of them beats a forecast that uses only the station and the calendar month. A month$\times$station climatology fitted on the training partition reaches AUC 0.685 and a Brier score of 0.2184. No learned model reaches it: Random Forest scores 0.007 lower (DeLong $p=0.57$), XGBoost, the LSTM, the Linear-4D control and \texttt{LSTM+QNN} 0.017--0.030 lower; only for the Linear-4D control does the year-block interval exclude zero (for \texttt{LSTM+QNN} it ends at zero).

On proper scores (\Cref{tab:brier}) every model, as trained, is worse than the climatology, and the year-block interval of the Brier difference excludes zero for all of them, Random Forest included (Brier skill $-0.046$; the others $-0.089$ to $-0.111$). All models were trained with positive-class weighting, which shifts probabilities away from the base rate, so part of this deficit reflects the training loss. After a Platt map fitted on the validation partition (1997--2007), the calibration errors fall from 0.10--0.14 to 0.03--0.06 and the Brier scores by 0.010--0.021, but no model beats the climatology on average, and at most one seed of any model does: Random Forest is at it ($-0.002$), the others below it ($-0.013$ to $-0.028$). The climatology itself is miscalibrated on the test years (ECE 0.061); recalibrated with the same map, its Brier score becomes 0.2155, and every model then has negative skill against it ($-0.016$ to $-0.042$).

\begin{table}[htbp]
\centering
\caption{Probabilistic skill against the month$\times$station climatology (15 seeds per model, 529 test rows). As trained: mean Brier score, Brier skill score (BSS), year-block 95\,\% interval of the Brier difference model $-$ climatology, 10-bin expected calibration error (ECE). After Platt recalibration on the validation partition: Brier score, BSS against the raw and against the recalibrated climatology ($n$: seeds with BSS $>0$), year-block interval of the Brier difference to the recalibrated climatology, ECE. Models of the recalibration columns refitted with the same seeds.}
\label{tab:brier}
{\small
\adjustbox{max width=\textwidth}{%
\begin{tabular}{@{}lccccccccc@{}}
\toprule
& \multicolumn{4}{c}{\textbf{As trained}} & \multicolumn{5}{c}{\textbf{After Platt recalibration}} \\
\cmidrule(lr){2-5}\cmidrule(lr){6-10}
\textbf{Model} & \textbf{Brier} & \textbf{BSS} & $\Delta$\textbf{BS [block CI]} & \textbf{ECE} & \textbf{Brier} & \textbf{BSS vs raw} ($n$) & \textbf{BSS vs recal.} ($n$) & $\Delta$\textbf{BS vs recal. [CI]} & \textbf{ECE} \\
\midrule
Climatology & 0.2184 & 0 & -- & 0.061 & 0.2155 & -- & 0 & -- & 0.042 \\
\midrule
Random Forest & 0.2284 & $-0.046$ & $[+0.002,+0.017]$ & 0.099 & 0.2189 & $-0.002$ (0) & $-0.016$ (0) & $[-0.003,+0.009]$ & 0.029 \\
XGBoost & 0.2394 & $-0.096$ & $[+0.011,+0.031]$ & 0.109 & 0.2212 & $-0.013$ (1) & $-0.027$ (0) & $[-0.000,+0.011]$ & 0.050 \\
Standalone LSTM & 0.2427 & $-0.111$ & $[+0.015,+0.033]$ & 0.144 & 0.2215 & $-0.014$ (1) & $-0.028$ (0) & $[+0.001,+0.010]$ & 0.053 \\
Linear-4D control & 0.2378 & $-0.089$ & $[+0.010,+0.028]$ & 0.125 & 0.2233 & $-0.023$ (0) & $-0.037$ (0) & $[+0.001,+0.014]$ & 0.055 \\
LSTM+QNN & 0.2424 & $-0.110$ & $[+0.015,+0.033]$ & 0.132 & 0.2245 & $-0.028$ (1) & $-0.042$ (0) & $[+0.003,+0.015]$ & 0.049 \\
\bottomrule
\end{tabular}}%
}
\end{table}

\paragraph{Within-station discrimination.} Because the threshold is pooled across stations whose base rates differ, the pooled AUC partly rewards telling stations apart. The within-station AUC counts only positive--negative pairs of months from the same station (\Cref{tab:within}). On this scale every learned model loses 0.020--0.032, the climatology 0.018, and the climatology has the highest value (0.667). Random Forest is not distinguished from it ($-0.020$, year-block $[-0.048,+0.011]$), whereas \texttt{LSTM+QNN} lies below it ($-0.041$, $[-0.080,-0.015]$).

\begin{table}[htbp]
\centering
\caption{Pooled and within-station AUC (15-seed means, 529 test rows). Within-station AUC: positive--negative pairs of the same station only, per-station AUCs averaged with weights $n_{\text{pos}}n_{\text{neg}}$. SD: across seeds. Per-station AUC in the last four columns (test $n$: 220, 221, 37, 51).}
\label{tab:within}
{\small
\adjustbox{max width=\textwidth}{%
\begin{tabular}{lcccccc}
\toprule
\textbf{Model} & \textbf{Pooled AUC} & \textbf{Within-station AUC (SD)} & \textbf{Alicante} & \textbf{Barcelona} & \textbf{San Javier} & \textbf{Valencia} \\
\midrule
Month$\times$station climatology & 0.685 & 0.667 (--) & 0.696 & 0.648 & 0.575 & 0.651 \\
\midrule
Random Forest & 0.678 & 0.647 (0.002) & 0.690 & 0.620 & 0.513 & 0.594 \\
XGBoost & 0.668 & 0.648 (0.014) & 0.682 & 0.625 & 0.529 & 0.623 \\
Standalone LSTM & 0.664 & 0.634 (0.010) & 0.653 & 0.622 & 0.564 & 0.621 \\
Linear-4D control & 0.655 & 0.623 (0.008) & 0.632 & 0.618 & 0.609 & 0.602 \\
LSTM+QNN & 0.659 & 0.627 (0.022) & 0.653 & 0.608 & 0.600 & 0.609 \\
LSTM+QTL & 0.655 & 0.627 (0.019) & 0.663 & 0.602 & 0.551 & 0.609 \\
\bottomrule
\end{tabular}}%
}
\end{table}

\subsection{Skill beyond the Seasonal and Station Base Rates}
\label{sec:stationmonth}

Pooling cases with different base rates credits a forecast with skill it does not have \citep{hamill2006}. Defining the target as exceedance of the training 60th percentile of the same station and calendar month removes that signal by construction. Under this label no model discriminates (\Cref{tab:stationmonth}): the 15-seed mean AUCs lie between 0.505 and 0.514, against 0.480 for the climatology, and no paired difference is distinguished from zero (DeLong $p\ge0.70$). After Platt recalibration every model has a lower Brier score than the raw climatology, but with AUCs near 0.5 this is a correction of the base rate towards that of the validation years; against the climatology recalibrated with the same map the skill is $+0.004$ to $+0.006$, and the year-block interval of the Brier difference includes zero for every model. The predictors therefore carry no detectable information about monthly exceedance beyond the seasonal cycle and the station, and the AUCs of 0.65--0.68 of \Cref{tab:case} are consistent with the models reproducing those base rates.

\begin{table}[htbp]
\centering
\caption{Station$\times$month label: 15-seed test AUC and Brier skill after Platt recalibration on the validation partition, against the raw and the recalibrated climatology ($n$: seeds with BSS $>0$), with the year-block interval of the Brier difference to the recalibrated climatology (529 test rows, 37.8\,\% positive).}
\label{tab:stationmonth}
{\small
\adjustbox{max width=\textwidth}{%
\begin{tabular}{lcccc}
\toprule
\textbf{Model} & \textbf{Mean AUC (SD)} & \textbf{BSS Platt vs raw clim.} ($n$) & \textbf{BSS Platt vs recal. clim.} ($n$) & $\Delta$\textbf{BS vs recal. [CI]} \\
\midrule
Month$\times$station climatology & 0.479 & -- & -- & -- \\
\midrule
Random Forest & 0.514 (0.005) & $+0.017$ (15) & $+0.006$ (15) & $[-0.004,+0.001]$ \\
XGBoost & 0.510 (0.013) & $+0.015$ (15) & $+0.004$ (12) & $[-0.004,+0.002]$ \\
Standalone LSTM & 0.510 (0.020) & $+0.015$ (15) & $+0.004$ (14) & $[-0.003,+0.002]$ \\
Linear-4D control & 0.512 (0.015) & $+0.015$ (15) & $+0.004$ (14) & $[-0.004,+0.003]$ \\
LSTM+QNN & 0.505 (0.018) & $+0.015$ (15) & $+0.004$ (13) & $[-0.003,+0.002]$ \\
\bottomrule
\end{tabular}}%
}
\end{table}

\subsection{An Operational Reference: ECMWF SEAS5}
\label{sec:seas5}

The climatology is the minimal reference; an agency would compare a statistical model with its dynamical seasonal forecast. We therefore add ECMWF SEAS5 \citep{johnson2019seas5}, distributed by the Copernicus Climate Change Service \citep{c3s_seasonal_monthly}, as a further model in the same protocol (\Cref{S-app:seas5}). Monthly precipitation of the first 25 members on the $1^\circ$ grid is interpolated bilinearly to each station. For lead 1 the forecast for month $M$ starts on the first day of $M$; it is released several days into $M$ (on the 5th to ECMWF Member States and on the 6th--13th through C3S, depending on the dataset and the period; \citealp{c3s_seasonal_monthly}), whereas the statistical models are available on the first day, so lead 1 is an upper bound on what the dynamical forecast offers at their issue time; lead 2 (start on the first day of $M-1$) is available before $M$ begins. The probability of exceedance, specified before the analysis, is the fraction of members above the model's own quantile at the level $1-p_{\text{clim}}$, with $p_{\text{clim}}$ the training exceedance frequency of the station and calendar month and the model quantile taken from the hindcasts of 1981--1996, so that no test or validation information enters the bias correction. Member fractions of a finite ensemble are not fair scores in the sense of \citet{ferro2014}; probabilistic scores use SEAS5 after Platt recalibration on the validation partition. To ask whether the statistical models add information to the dynamical forecast, we fit for each model and seed, on the validation partition, a logistic combination of the logits of the model and of SEAS5, and compare it with the same combination of the climatology and SEAS5.

\begin{table}[htbp]
\centering
\caption{SEAS5 as an operational reference (529 test rows; lead 1 unless stated). Blend: logistic combination of the logits of the named model and of SEAS5, fitted on the validation partition, one per seed. $\hat\Delta$: difference of 15-seed mean AUC (SEAS5 and the climatology are single scorers). Cluster/iid: ratio of the cluster-robust to the DeLong standard error. Secondary analysis: $p$-values unadjusted.}
\label{tab:seas5}
{\small
\adjustbox{max width=\textwidth}{%
\begin{tabular}{@{}lccccccc@{}}
\toprule
Comparison & $\hat\Delta$ & DeLong $p$ & DeLong 95\,\% CI & Year-block 95\,\% CI & Years$\times$seeds 95\,\% CI & Cluster/iid & MDE \\
\midrule
SEAS5 $-$ climatology & $+0.0433$ & 0.034 & $[+0.003,+0.083]$ & $[-0.005,+0.086]$ & $[-0.004,+0.085]$ & 1.17 & 0.057 \\
SEAS5 (lead 2) $-$ climatology & $-0.0095$ & 0.460 & $[-0.035,+0.016]$ & $[-0.034,+0.015]$ & $[-0.034,+0.017]$ & 1.05 & 0.036 \\
Random Forest $-$ SEAS5 & $-0.0501$ & 0.021 & $[-0.093,-0.007]$ & $[-0.094,-0.006]$ & $[-0.095,-0.006]$ & 1.11 & 0.061 \\
LSTM+QNN $-$ SEAS5 & $-0.0698$ & 0.001 & $[-0.112,-0.028]$ & $[-0.116,-0.022]$ & $[-0.117,-0.020]$ & 1.17 & 0.060 \\
\midrule
Blend RF $-$ blend climatology & $+0.0006$ & 0.877 & $[-0.007,+0.009]$ & $[-0.008,+0.008]$ & $[-0.008,+0.009]$ & 1.02 & 0.012 \\
Blend XGBoost $-$ blend climatology & $-0.0002$ & 0.970 & $[-0.009,+0.009]$ & $[-0.012,+0.011]$ & $[-0.012,+0.011]$ & 1.27 & 0.013 \\
Blend LSTM $-$ blend climatology & $+0.0007$ & 0.831 & $[-0.006,+0.008]$ & $[-0.006,+0.008]$ & $[-0.005,+0.007]$ & 0.87 & 0.010 \\
Blend Linear-4D $-$ blend climatology & $+0.0008$ & 0.858 & $[-0.008,+0.009]$ & $[-0.008,+0.010]$ & $[-0.008,+0.010]$ & 1.02 & 0.012 \\
Blend LSTM+QNN $-$ blend climatology & $+0.0002$ & 0.952 & $[-0.006,+0.007]$ & $[-0.006,+0.008]$ & $[-0.007,+0.008]$ & 1.00 & 0.010 \\
Blend LSTM+QNN $-$ blend Linear-4D & $-0.0006$ & 0.744 & $[-0.004,+0.003]$ & $[-0.005,+0.004]$ & $[-0.006,+0.004]$ & 1.29 & 0.005 \\
\bottomrule
\end{tabular}}}
\end{table}

At lead 1 SEAS5 scores AUC 0.728, against 0.685 for the climatology and 0.678 for Random Forest (\Cref{tab:seas5}). Its gain over the climatology, $+0.043$, is below its MDE (0.057) and is not resolved: DeLong rejects the null ($p=0.034$), but the year-block and two-way intervals include zero, the cluster-robust standard error is 17\,\% above the DeLong one and the cluster-robust interval includes zero for every start month of the 12-month clusters. Random Forest and \texttt{LSTM+QNN}, in turn, lie 0.050 and 0.070 below SEAS5; the second difference is resolved by every procedure, the first by DeLong and both bootstrap intervals and by the cluster-robust interval for 7 of the 12 start months. The dynamical forecast is the only predictor with detectable information beyond the base rates: with the station$\times$month label it scores 0.639, against 0.480 for the climatology ($+0.159$, year-block $[+0.087,+0.226]$) and 0.505--0.514 for every statistical model. The station predictors add nothing to it: combined with SEAS5, every statistical model lies within $-0.0002$ to $+0.0008$ of the combination of the climatology with SEAS5 (MDEs 0.010--0.013), and the combinations with \texttt{LSTM+QNN} and with the Linear-4D control differ by $-0.0006$ (MDE 0.005). After recalibration SEAS5 has a Brier skill of $+0.051$ against the climatology (year-block interval of the Brier difference $[-0.024,+0.003]$). At lead 2 the advantage disappears ($-0.010$, $[-0.034,+0.015]$; Brier skill $-0.016$). An uncorrected count of members above the threshold scores higher at lead 1 (0.748; \Cref{S-tab:seas5_full}); it was not the specified variant and is reported only descriptively. For the benchmark the result has two sides: the station models do not reach the operational forecast, and even the operational forecast, at the most favourable lead, is not resolvably better than the climatology with 529 test months. The climatology remains the reference that a statistical model at this lead would first have to beat.

\subsection{Out-of-Panel Re-test}
\label{sec:replication}

We repeated the protocol on a second Western Mediterranean panel, Palma de Mallorca airport, Tortosa, Alcantarilla and Almería airport (AEMET daily series from 1942, 1961, 1968 and 1972 to May 2026), with the same code path, labels, split years, six models and the Random Forest and XGBoost configurations tuned on the first panel (\Cref{S-app:panel2}). All four stations cover the test period, so the test set has 884 station-months (34.7\,\% positive) and the shorter histories leave 1{,}694 training rows. Before the analysis we recorded, and committed to the repository, the resolution the protocol predicts for this test set (\path{prediction_panel2.txt}, 10 November 2026): a standard error of one AUC of about 0.019 (Hanley--McNeil at AUC 0.70) and MDEs of 0.025--0.035 for paired 15-seed means of the classical models against each other and the neural models, by $N^{-1/2}$ scaling of the first panel.

\begin{table}[htbp]
\centering
\caption{Second panel (Mallorca, Tortosa, Alcantarilla, Almería; 884 test rows, 19 years, 15 seeds; pooled P60 label). $\hat\Delta$: difference of 15-seed mean AUC, with the first-panel value. Cluster/iid: calendar-year cluster-robust to DeLong standard error; phases: start months (of 12) for which the cluster-robust $t_{18}$ interval excludes zero. Secondary analysis: $p$-values unadjusted.}
\label{tab:panel2}
{\small
\adjustbox{max width=\textwidth}{%
\begin{tabular}{@{}lcccccccc@{}}
\toprule
Comparison & $\hat\Delta$ & First panel & DeLong $p$ & Year-block 95\,\% CI & Years$\times$seeds 95\,\% CI & MDE & Cluster/iid & Phases \\
\midrule
RF $-$ XGBoost & $+0.0041$ & $+0.0100$ & 0.449 & $[-0.009,+0.017]$ & $[-0.010,+0.018]$ & 0.015 & 1.33 & 0 \\
RF $-$ LSTM & $-0.0101$ & $+0.0146$ & 0.278 & $[-0.027,+0.006]$ & $[-0.028,+0.006]$ & 0.026 & 0.91 & 0 \\
RF $-$ LSTM+QNN & $+0.0008$ & $+0.0197$ & 0.933 & $[-0.020,+0.020]$ & $[-0.021,+0.020]$ & 0.027 & 1.09 & 0 \\
XGBoost $-$ LSTM+QNN & $-0.0033$ & $+0.0097$ & 0.763 & $[-0.026,+0.019]$ & $[-0.026,+0.021]$ & 0.030 & 1.11 & 0 \\
LSTM $-$ LSTM+QNN & $+0.0110$ & $+0.0051$ & 0.022 & $[-0.002,+0.023]$ & $[-0.003,+0.025]$ & 0.013 & 1.35 & 0 \\
Linear-4D $-$ LSTM+QNN & $+0.0068$ & $-0.0034$ & 0.234 & $[-0.007,+0.020]$ & $[-0.008,+0.021]$ & 0.016 & 1.19 & 0 \\
Fourier-2 $-$ LSTM+QNN & $+0.0031$ & $+0.0041$ & 0.259 & $[-0.001,+0.007]$ & $[-0.006,+0.011]$ & 0.008 & 0.74 & 0 \\
\bottomrule
\end{tabular}}}
\end{table}

The predicted resolution held: three of the four MDEs of the classical comparisons lie in the predicted range (0.026--0.030), and Random Forest against XGBoost, the most correlated pair after tuning, lies below it (0.015; \Cref{tab:panel2}). So did the null results: the climatology scores 0.716 and only the standalone LSTM has a higher mean (0.717); with the station$\times$month label every model scores 0.499--0.523 and no paired comparison excludes zero (\Cref{S-tab:panel2_sm}); and the hybrid is not distinguished from the Linear-4D control or from \texttt{Fourier-2} by any procedure. What did not replicate is the ranking: the six models lie within 0.014 of one another, Random Forest is no longer above the hybrid ($+0.001$) and the LSTM is above Random Forest ($+0.010$), all inside the MDE. The standalone LSTM exceeds the hybrid by 0.011 (DeLong $p=0.022$; year-block interval reaching zero; not resolved). The cluster/iid ratios are 0.91--1.35, so the procedural disagreement of the first panel does not recur. The second panel shares the region and the 2008--2026 test years with the first, so its sampling errors are not independent of the first panel's and it is a re-test, not an independent replication; it shows, as the protocol predicts, that differences below the MDE of one panel need not keep their sign on another.

% ============================================================
\section{Discussion}
% ============================================================

\subsection{What a Benchmark of This Size Can and Cannot Conclude}

With 529 correlated test months the 80\,\%-power MDE for 15-seed means is 0.022--0.032 AUC, and the observed gaps between the learned models (at most 0.024) are mostly below it. This does not mean that the models are equal. Three observations follow. First, on the scale of one AUC, test-sampling variability exceeds seed variability for every model, but for the hybrids the seed share rises to 14--25\,\%, so seed and sampling error must both be reported. Second, near the MDE the verdict depends on the estimand and on the procedure: single-seed verdicts depend on which seeds are compared, and for the seed means DeLong and the cluster procedures disagree on Random Forest against the LSTM and the hybrids. The disagreement comes from a cross-station co-movement within years that the iid estimator ignores and that did not recur on the second panel, so the choice of procedure is a substantive assumption about the dependence between test cases, and an agreement of the procedures is better evidence than either one. Third, the references are strong: no learned model improves on the month$\times$station climatology in AUC or Brier score, and with a label that removes the base rates none discriminates.

For the quantum results, the hybrids lie 0.020--0.024 AUC below Random Forest and are not distinguished from circuit-free classical controls at the same interface: Fourier features of the four angles match the trained circuit, with or without cross-coordinate products, and linear and tensor heads on a frozen extractor match the quantum head, all within the resolution of the comparison. Training the circuit helps relative to freezing it. A planted signal shows that the controls would have separated a circuit contribution from a linear or separable map had there been one, and that a tensor surrogate reproduces it. The data are therefore consistent with a trained periodic map at the interface being what helps, with no detectable contribution from its being quantum. For a 4-qubit circuit this is expected from the classical-surrogate results \citep{schreiber2023,sweke2025}; the point of the controls is that the benchmark shows it directly, and that the same controls would be needed, and harder to replace by an argument, for circuits whose spectrum or entanglement makes classical surrogates expensive. The evidence does not establish a quantum performance advantage in this case study; it neither shows nor rules out one on larger or richer data.

\subsection{What These Results Do Not Show}
\label{sec:not_show}

Because Random Forest has the highest mean AUC among the learned models, the results could be read as one more instance of tree ensembles outperforming other models on tabular data \citep{grinsztajn2022}. They do not support that reading. Random Forest does not reach the month$\times$station climatology in AUC, it is worse on Brier score as trained and not better after recalibration, and within a station it discriminates no better than the climatology (\Cref{tab:within}). Its advantage over the LSTM and \texttt{LSTM+QNN} is resolved by the cluster procedures but not by DeLong and does not survive the Holm correction, and on the second panel it vanishes. What the data support is that none of the learned models, classical or quantum, adds resolvable skill beyond the seasonal and geographic base rate of this panel, and that the differences between them are below what 529 test months can resolve.

Nor do the results show that variational circuits are ineffective in general: the study concerns 4-qubit circuits behind a 4-dimensional classical bottleneck, simulated without noise, on a small panel whose predictable signal is largely seasonal. And they do not show that DeLong or the cluster procedures are wrong in general: each is calibrated under some of the simulated mechanisms and not under others.

\subsection{A Low-Signal Testbed}
\label{sec:testbed}

The strongest empirical result of the case study is negative for every statistical model: none improves on the climatology, none discriminates once the base rates are removed, and none adds information to the dynamical forecast. At a one-month lead, with monthly aggregates from four stations and these predictors, the signal available beyond the seasonal and geographic base rate is too small to be detected with this test set; even SEAS5, started on the first day of the target month, is not resolvably better than the climatology on the pooled label. The tuning of the tree ensembles points the same way: the search selected strongly regularized forests (at least 10 rows per leaf), and early stopping on the validation AUC selected very few boosting rounds. Several factors may contribute: the limited predictability of monthly precipitation one month ahead from local antecedent data, and monthly sums that do not resolve the convective episodes that produce much of the heavy precipitation in this region.

This makes the panel a poor testbed for deciding which architecture is better: on these data the answer is none of them. For the question of this paper it is a relevant case rather than a weakness. Small QML benchmarks, including those cited in the Introduction, often compare models whose differences are of this size, and the noise budget, the disagreement between procedures and the resolution rule describe what such a benchmark can conclude whatever the reason for the small differences. The statistical argument does not depend on the panel having little signal; the conclusions about particular architectures do, and they should be tested on data where learned models clearly beat a base-rate reference.

\subsection{Implications for Benchmark Design}
\label{sec:design}

The resolution limits above follow from the size and composition of the test set, not from the climate application, so they can be estimated before a benchmark is run. For a binary classifier with AUC $A$, $n_1$ positives and $n_0$ negatives, the Hanley--McNeil approximation \citep{hanley1982} gives the standard error of one AUC from $(A,n_1,n_0)$ alone; for the present test set it returns 0.024, matching the DeLong value. For two models scored on the same rows, the standard error of the difference follows from the two standard errors and the correlation between the AUCs \citep{hanley1983}. Assuming a correlation of 0.75, \Cref{tab:hm_mde} gives the approximate 80\,\%-power MDE for other test-set sizes: about 0.049 for $N=500$ at prevalence 0.41 and about 0.086 at prevalence 0.10. These values are specific to the assumed AUC (0.68), inter-model correlation and iid approximation; temporal and spatial dependence can change them in either direction for paired comparisons.

\begin{table}[htbp]
\centering
\caption{Approximate MDE for the AUC difference of two models on the same test set (Hanley--McNeil standard errors combined with an inter-model correlation of 0.75 as in \citealp{hanley1983}; AUC 0.68, $\alpha=0.05$, power 0.80). Test-set size $N$ and positive-class prevalence $\pi$.}
\label{tab:hm_mde}
{\small
\begin{tabular}{rcc}
\toprule
$N$ & MDE, $\pi=0.41$ & MDE, $\pi=0.10$ \\
\midrule
250 & 0.069 & 0.122 \\
500 & 0.049 & 0.086 \\
1{,}000 & 0.035 & 0.061 \\
2{,}500 & 0.022 & 0.038 \\
5{,}000 & 0.015 & 0.027 \\
\bottomrule
\end{tabular}%
}
\end{table}

An empirical version of this scaling is obtained from the panel itself (\Cref{tab:resolution}): random subsets of 5 to 16 of the 19 test years (all stations of a year together) are re-evaluated with the archived predictions. The MDEs fall approximately as $N^{-1/2}$, and the single-run MDEs lie 2--26\,\% below the Hanley--McNeil values, so \Cref{tab:hm_mde} is slightly conservative here. Sample sizes should be planned for a smallest difference of interest fixed in advance, not for the observed gaps. Scaling the full-test-set standard errors as $N^{-1/2}$, 80\,\% power for a true Random Forest $-$ \texttt{LSTM+QNN} difference of 0.02 needs about 1{,}370 test rows under DeLong (2.6 times the present test set) and about 720 under the cluster-robust standard error; a true difference of 0.01 between the hybrid and the Linear-4D control needs about 1{,}430 rows under DeLong and 1{,}700 under the cluster-robust standard error, and one of 0.005 four times as many. These extrapolations assume that the correlation between models and the dependence between years stay as on the panel.

\begin{table}[htbp]
\centering
\caption{Empirical resolution curve: median 80\,\%-power DeLong MDE over random subsets of test years (60 subsets per size; archived predictions). Single run: seed $s$ against seed $s$, averaged over seeds. 15-seed means: paired DeLong on seed-averaged components. The last row is the full test set; HM: Hanley--McNeil single-run value at the same $N$ (prevalence 0.406, AUC 0.68, correlation 0.75).}
\label{tab:resolution}
{\small
\begin{tabular}{rrcccc}
\toprule
& & & \multicolumn{2}{c}{\textbf{Single run}} & \textbf{15-seed means} \\
\cmidrule(lr){4-5}\cmidrule(lr){6-6}
\textbf{Years} & $N$ & \textbf{HM} & \textbf{RF $-$ LSTM} & \textbf{RF $-$ LSTM+QNN} & \textbf{RF $-$ LSTM+QNN} \\
\midrule
5 & 143 & 0.092 & 0.069 & 0.090 & 0.063 \\
8 & 219 & 0.074 & 0.055 & 0.071 & 0.050 \\
10 & 278 & 0.066 & 0.049 & 0.063 & 0.044 \\
13 & 361 & 0.058 & 0.042 & 0.055 & 0.039 \\
16 & 454 & 0.051 & 0.038 & 0.050 & 0.035 \\
19 & 529 & 0.048 & 0.035 & 0.046 & 0.032 \\
\bottomrule
\end{tabular}%
}
\end{table}

In practice, the MDE of a planned comparison can be computed from the test size, the prevalence and an assumed AUC before any model is trained (\Cref{tab:hm_mde}) and checked afterwards on the archived predictions (\Cref{tab:resolution}). If the differences of interest are below it, more seeds do not help: averaging seeds reduces the seed component but leaves the sampling error unchanged, and a larger test set is the only remedy.

\subsection{Recommended Reporting Protocol}
\label{sec:protocol}

\Cref{tab:protocol} lists the practices that follow from the results. Items 1, 2 and 4 need only the test probabilities of the models; item 3 needs a specified and executed generating mechanism.

\begin{table}[htbp]
\centering
\caption{Reporting protocol for small QML benchmarks.}
\label{tab:protocol}
{\small
\begin{tabularx}{\textwidth}{cX}
\toprule
\textbf{\#} & \textbf{Practice} \\
\midrule
1 & State the estimand (single run or seed average; conditional on which sources) and report paired confidence intervals for it, with DeLong and a cluster procedure over the natural cluster (here calendar years); report discrepancies between procedures instead of selecting one, and diagnose them, for example by splitting the cluster-robust variance into within- and cross-site terms. \\
2 & Report the MDE as a measure of design sensitivity next to the intervals; do not treat it as a significance threshold or classify gaps below it as uninformative, and plan test sizes for a smallest difference of interest fixed in advance. \\
3 & Check the procedures against a simulation with known differences at the sample size and dependence structure of the study. \\
4 & Report across-seed variability separately from test-sampling uncertainty; treat seed-level tests as statements about optimization, not generalization. \\
5 & Include a climatological (or base-rate) reference with a proper scoring rule, and report skill relative to it after recalibrating both the models and the reference on held-out data; where the label has strong base rates, also report skill with a label that removes them. Where an operational dynamical forecast exists, include it as a model and test whether the statistical model adds to it. \\
6 & Fit all model families on the same rows and tune them with a declared search of comparable budget, or state and quantify the asymmetry; run all models in one pinned software environment. \\
7 & Add classical controls at the same low-dimensional interface as the quantum circuit: one with a matched parameter count, Fourier surrogates with the spectrum of the circuit including its cross-coordinate products, and the circuit with frozen weights; check with a planted circuit signal that the controls can detect a circuit contribution; attribute differences to the circuit only through these controls. \\
8 & Build labels from the target variable alone, with an explicit completeness rule, and keep the predictors' missing data out of the label. \\
9 & Choose equivalence margins by scientific relevance, independently of the MDE and of the results, and treat an interval wider than the margin as inconclusive. \\
10 & When the target threshold is pooled across heterogeneous sites, report a within-site AUC, so that skill is not credited for separating sites with different base rates. \\
\bottomrule
\end{tabularx}
}
\end{table}

\subsection{Limitations}
\label{sec:limitations}

\textbf{Inferential procedures.} The cluster procedures rely on 19 calendar years. Their disagreement with DeLong is traced on the first panel to a co-movement of the station contributions within years, which is at the edge of what a permutation null produces and did not recur on the second panel; a test period with more independent years, or several independent test periods, would be needed to settle the RF $-$ LSTM and RF $-$ \texttt{LSTM+QNN} verdicts.

\textbf{Hybrid models.} The comparisons concern 4-qubit circuits behind a 4-dimensional classical bottleneck in noiseless simulation; the hybrids have about a quarter of the trainable parameters of the standalone LSTM, so only the comparisons with the interface controls are capacity-matched. The controls were added after the main analysis, and no equivalence margin was specified in advance. The tensor-product surrogates of \texttt{LSTM+QNN} contain the cross-coordinate products only up to degree 2 per coordinate, while the circuit's spectrum extends to degree 4; the planted-signal control shows that on this encoding the low-degree surrogates reproduce a circuit-generated signal, but not that they would for other encodings. The planted signal was recovered at the interface and not by complete models trained on the raw inputs.

\textbf{Data and labels.} Two of the four stations of the first panel have no complete predictors after 2011--2012, so 83\,\% of the test windows come from two stations; filling the pressure records from neighbouring stations or reanalysis was not attempted. Months with more than 3 missing precipitation days have no label, which removes a small, non-random set of months.

\textbf{Baselines and scope.} The tree ensembles were tuned on the validation partition, which XGBoost also uses for early stopping; the search reached a validation AUC of 0.701 (Random Forest) and 0.723 (XGBoost) and moved the test AUC by $+0.001$ and $+0.015$ relative to the untuned configurations, neither difference resolved (\Cref{S-app:hyperparams}). The neural models were not tuned automatically. All models except the climatology were trained with positive-class weighting; the Brier comparisons describe the models as trained and after Platt recalibration. The panel is small, monthly aggregation does not resolve convective episodes, and the chronological split falls in a period of accelerated warming. We evaluate one ansatz family in noiseless simulation; noise during training \citep{wang2021noise} and hardware execution were not studied. The SEAS5 reference uses one dynamical system, 25 members, a $1^\circ$ grid and a bias correction calibrated on 16 hindcast years.

\textbf{Generality.} The empirical conclusions rest on two panels of four Western Mediterranean stations each, which share the split years and the test period. Other climates, targets and test periods remain untested. The sampling argument and the resolution rule transfer to other benchmarks (\Cref{sec:design}), but the findings about particular models and references do not replace validation on further datasets, preferably ones in which learned models clearly improve on a base-rate reference.

% ============================================================
\section{Conclusions}
% ============================================================

On a benchmark of 529 correlated test station-months, the test sample, not the training seed, limits what can be concluded about the classical models: one AUC has a standard error of about 0.024, the minimum detectable difference of 15-seed means is 0.022--0.032 AUC, and it falls only as $N^{-1/2}$ with the size of the test set. For variational hybrids the seed variance is also substantial, so both sources must be reported. Near this limit the verdict depends on the procedure: DeLong and the calendar-year cluster procedures disagree on Random Forest against the LSTM and the hybrids, because of a co-movement of stations within years that the iid estimator ignores and that did not recur on a second panel.

In the quantum case study, two hybrid models lie 0.020--0.024 AUC below Random Forest and are not distinguished from circuit-free classical controls at the same interface; Fourier surrogates on the circuit's angles match the trained circuit, whose readouts are nearly degree-1 tensor polynomials of its angles whatever its weights, and the same controls detect a circuit-generated signal when one is planted. Finite-shot inference costs 0.002 AUC at 1{,}000 shots. No learned model improves on a month$\times$station climatology in AUC or Brier score; with a station- and month-specific threshold none discriminates; and none adds information to the ECMWF SEAS5 forecast, which itself is not resolvably better than the climatology on the pooled label with this test set. On a second panel the predicted resolution and the null results held, and the model ranking did not.

The consequence for QML benchmarking is practical: state the estimand, compute the minimum detectable difference before ranking architectures, report seed and sampling error separately, compare against base-rate and operational references and against classical controls at the same interface, and treat differences below the resolution of the design as unresolved rather than as evidence for either model.

\section*{Data and Code Availability}
\label{sec:data_availability}
Code, the processed monthly panel and the archived test predictions of every model and seed are available at \url{https://github.com/enriquerb32/limits-of-resolution-qml}\todo{archive the release of this repository with a DOI (for example Zenodo) and cite that identifier}. The results of this paper come from the \texttt{v2} pipeline (\path{code/etl_v2.py}, \path{code/local_experiments/dataset_v2.py}), run end to end by \path{run_v2.ps1} in the environment fixed by \path{requirements-v2.lock}; the results of the earlier version are kept in \path{results_v1/}. The daily records are AEMET open data obtained through \textit{datosclima.es}; only aggregated monthly data are redistributed, with attribution to AEMET. The daily series of the second panel (AEMET stations B278, 9981A, 7228 and 6325O) were obtained in the same way. SEAS5 data were obtained from the Copernicus Climate Change Service Climate Data Store \citep{c3s_seasonal_monthly} under the Licence to Use Copernicus Products; the station-interpolated members are archived in the repository. The repository README maps every table and figure to its command and output file, and \path{scripts/build_traceability_v2.py} checks the numerical claims of the paper against the archived results (after \path{exp42_v2_summary.py} and \path{exp43_v2_tables.py}, which compute the summary statistics from the archived predictions).

\section*{Declaration of Generative AI in Scientific Writing}
During the preparation of this manuscript, the author used generative AI tooling for editorial language polishing, LaTeX formatting and structural drafting, and for drafting analysis scripts that the author ran and checked. No reported test statistic, empirical metric or benchmark result was generated by AI models; the numerical results come from the released scripts. The author assumes full scientific and editorial responsibility for the content.

\section*{Author Contributions}
E.R. conceived the study, curated the data, developed the models, performed the statistical analysis, and wrote the manuscript.

\section*{Competing Interests}
The author is a trainee meteorologist of the Spanish State Meteorological Agency (AEMET) and works as a data engineer at Minsait. This research was carried out independently of both positions and does not represent the views of either organization. The author declares no other competing financial or non-financial interests.

\section*{Funding}
This research received no external funding.

\section*{Acknowledgements}
The author acknowledges the State Meteorological Agency of Spain (AEMET) for providing the underlying meteorological station records under open-data terms.

\newpage
\begin{thebibliography}{99}

\bibitem[Alexandersson(1986)]{alexandersson1986}
Alexandersson, H. (1986). A homogeneity test applied to precipitation data. \textit{International Journal of Climatology}, 6(6), 661--675. \url{https://doi.org/10.1002/joc.3370060607}

\bibitem[Allen et al.(1998)]{allen1998}
Allen, R. G., Pereira, L. S., Raes, D., \& Smith, M. (1998). \textit{Crop evapotranspiration---Guidelines for computing crop water requirements---FAO Irrigation and drainage paper 56}. FAO, Rome.

\bibitem[Bates et al.(2024)]{bates2024}
Bates, S., Hastie, T., \& Tibshirani, R. (2024). Cross-validation: What does it estimate and how well does it do it? \textit{Journal of the American Statistical Association}, 119(546), 1434--1445. \url{https://doi.org/10.1080/01621459.2023.2197686}

\bibitem[Bazgir \& Zhang(2024)]{bazgir2024}
Bazgir, A., \& Zhang, Y. (2024). QESM: A leap towards quantum-enhanced ML emulation framework for Earth and climate modeling. \textit{arXiv preprint arXiv:2410.01551}. \url{https://arxiv.org/abs/2410.01551}

\bibitem[Bell \& McCaffrey(2002)]{bell2002}
Bell, R. M., \& McCaffrey, D. F. (2002). Bias reduction in standard errors for linear regression with multi-stage samples. \textit{Survey Methodology}, 28(2), 169--181.

\bibitem[Bouthillier et al.(2021)]{bouthillier2021}
Bouthillier, X., Delaunay, P., Bronzi, M., Trofimov, A., Nichyporuk, B., Szeto, J., Sepah, N., Raff, E., Madan, K., Voleti, V., Ebrahimi Kahou, S., Michalski, V., Serdyuk, D., Arbel, T., Pal, C., Varoquaux, G., \& Vincent, P. (2021). Accounting for variance in machine learning benchmarks. \textit{Proceedings of Machine Learning and Systems}, 3, 747--769.

\bibitem[Bowles et al.(2024)]{bowles2024}
Bowles, J., Ahmed, S., \& Schuld, M. (2024). Better than classical? The subtle art of benchmarking quantum machine learning models. \textit{arXiv preprint arXiv:2403.07059}. \url{https://arxiv.org/abs/2403.07059}

\bibitem[Breiman(2001)]{breiman2001}
Breiman, L. (2001). Random forests. \textit{Machine Learning}, 45(1), 5--32. \url{https://doi.org/10.1023/A:1010933404324}

\bibitem[Brier(1950)]{brier1950}
Brier, G. W. (1950). Verification of forecasts expressed in terms of probability. \textit{Monthly Weather Review}, 78(1), 1--3. \url{https://doi.org/10.1175/1520-0493(1950)078<0001:VOFEIT>2.0.CO;2}

\bibitem[Cameron et al.(2008)]{cameron2008}
Cameron, A. C., Gelbach, J. B., \& Miller, D. L. (2008). Bootstrap-based improvements for inference with clustered errors. \textit{The Review of Economics and Statistics}, 90(3), 414--427. \url{https://doi.org/10.1162/rest.90.3.414}

\bibitem[Copernicus Climate Change Service(2018)]{c3s_seasonal_monthly}
Copernicus Climate Change Service, Climate Data Store (2018). Seasonal forecast monthly statistics on single levels. Copernicus Climate Change Service (C3S) Climate Data Store (CDS). \url{https://doi.org/10.24381/cds.68dd14c3}\todo{check the dataset DOI on the CDS page and add the access date}

\bibitem[Card et al.(2020)]{card2020}
Card, D., Henderson, P., Khandelwal, U., Jia, R., Mahowald, K., \& Jurafsky, D. (2020). With little power comes great responsibility. In \textit{Proceedings of the 2020 Conference on Empirical Methods in Natural Language Processing (EMNLP)} (pp. 9263--9274). \url{https://doi.org/10.18653/v1/2020.emnlp-main.745}

\bibitem[Cerezo et al.(2025)]{cerezo2025}
Cerezo, M., Larocca, M., Garc\'ia-Mart\'in, D., Diaz, N. L., Braccia, P., Fontana, E., Rudolph, M. S., Bermejo, P., Ijaz, A., Thanasilp, S., Anschuetz, E. R., \& Holmes, Z. (2025). Does provable absence of barren plateaus imply classical simulability? \textit{Nature Communications}, 16, 7907. \url{https://doi.org/10.1038/s41467-025-63099-6}

\bibitem[da Silva et al.(2026)]{dasilva2026}
da Silva, M. H. F., de Jesus, G. F., Nascimento, C. M. S., da Silva, V. L., \& Cruz, C. S. (2026). Exploring quantum machine learning for weather forecasting. \textit{Brazilian Journal of Physics}, 56, 22. \url{https://doi.org/10.1007/s13538-025-01941-4}

\bibitem[DeLong et al.(1988)]{delong1988}
DeLong, E. R., DeLong, D. M., \& Clarke-Pearson, D. L. (1988). Comparing the areas under two or more correlated receiver operating characteristic curves: A nonparametric approach. \textit{Biometrics}, 44(3), 837--845. \url{https://doi.org/10.2307/2531595}

\bibitem[DelSole \& Tippett(2014)]{delsole2014}
DelSole, T., \& Tippett, M. K. (2014). Comparing forecast skill. \textit{Monthly Weather Review}, 142(12), 4658--4678. \url{https://doi.org/10.1175/MWR-D-14-00045.1}

\bibitem[Diebold \& Mariano(1995)]{diebold1995}
Diebold, F. X., \& Mariano, R. S. (1995). Comparing predictive accuracy. \textit{Journal of Business \& Economic Statistics}, 13(3), 253--263. \url{https://doi.org/10.1080/07350015.1995.10524599}

\bibitem[Dietterich(1998)]{dietterich1998}
Dietterich, T. G. (1998). Approximate statistical tests for comparing supervised classification learning algorithms. \textit{Neural Computation}, 10(7), 1895--1923. \url{https://doi.org/10.1162/089976698300017197}

\bibitem[Ferro(2014)]{ferro2014}
Ferro, C. A. T. (2014). Fair scores for ensemble forecasts. \textit{Quarterly Journal of the Royal Meteorological Society}, 140(683), 1917--1923. \url{https://doi.org/10.1002/qj.2270}

\bibitem[Field \& Welsh(2007)]{field2007}
Field, C. A., \& Welsh, A. H. (2007). Bootstrapping clustered data. \textit{Journal of the Royal Statistical Society: Series B (Statistical Methodology)}, 69(3), 369--390. \url{https://doi.org/10.1111/j.1467-9868.2007.00593.x}

\bibitem[Fellner et al.(2026)]{fellner2026}
Fellner, T., Kreplin, D. A., Tovey, S., \& Holm, C. (2026). Quantum vs. classical: A comprehensive benchmark study for predicting time series with variational quantum machine learning. \textit{Machine Learning: Science and Technology}, 7(1), 010501. \url{https://doi.org/10.1088/2632-2153/ae365f}

\bibitem[Gil-Fuster et al.(2025)]{gilfuster2025}
Gil-Fuster, E., Gyurik, C., P\'erez-Salinas, A., \& Dunjko, V. (2025). On the relation between trainability and dequantization of variational quantum learning models. In \textit{International Conference on Learning Representations (ICLR 2025)}. \url{https://arxiv.org/abs/2406.07072}

\bibitem[Grinsztajn et al.(2022)]{grinsztajn2022}
Grinsztajn, L., Oyallon, E., \& Varoquaux, G. (2022). Why do tree-based models still outperform
deep learning on typical tabular data? \textit{Advances in Neural Information Processing
Systems (NeurIPS)}, 35, 507--520.

\bibitem[Hamill \& Juras(2006)]{hamill2006}
Hamill, T. M., \& Juras, J. (2006). Measuring forecast skill: Is it real skill or is it the varying climatology? \textit{Quarterly Journal of the Royal Meteorological Society}, 132(621C), 2905--2923. \url{https://doi.org/10.1256/qj.06.25}

\bibitem[Hanley \& McNeil(1982)]{hanley1982}
Hanley, J. A., \& McNeil, B. J. (1982). The meaning and use of the area under a receiver operating characteristic (ROC) curve. \textit{Radiology}, 143(1), 29--36. \url{https://doi.org/10.1148/radiology.143.1.7063747}

\bibitem[Hanley \& McNeil(1983)]{hanley1983}
Hanley, J. A., \& McNeil, B. J. (1983). A method of comparing the areas under receiver operating characteristic curves derived from the same cases. \textit{Radiology}, 148(3), 839--843. \url{https://doi.org/10.1148/radiology.148.3.6878708}

\bibitem[Henderson et al.(2018)]{henderson2018}
Henderson, P., Islam, R., Bachman, P., Pineau, J., Precup, D., \& Meger, D. (2018). Deep reinforcement learning that matters. In \textit{Proceedings of the AAAI Conference on Artificial Intelligence}, 32(1), 3207--3214. \url{https://doi.org/10.1609/aaai.v32i1.12007}

\bibitem[Huang et al.(2021)]{huang2021}
Huang, H.-Y., Broughton, M., Cotler, J., Chen, S., Li, J., Mohseni, M., Neven, H., Babbush, R., Kueng, R., Preskill, J., \& McClean, J. R. (2021). Power of data in quantum machine learning. \textit{Nature Communications}, 12, 2631. \url{https://doi.org/10.1038/s41467-021-22539-9}

\bibitem[Johnson et al.(2019)]{johnson2019seas5}
Johnson, S. J., Stockdale, T. N., Ferranti, L., Balmaseda, M. A., Molteni, F., Magnusson, L., Tietsche, S., Decremer, D., Weisheimer, A., Balsamo, G., Keeley, S. P. E., Mogensen, K., Zuo, H., \& Monge-Sanz, B. M. (2019). SEAS5: the new ECMWF seasonal forecast system. \textit{Geoscientific Model Development}, 12(3), 1087--1117. \url{https://doi.org/10.5194/gmd-12-1087-2019}

\bibitem[Kumar et al.(2026)]{kumar2026}
Kumar, T. P., Raghavender, K. V., Veeranna, T., Srikanth, G. K., \& Routray, A. (2026). Hybrid variational quantum-classical data assimilation for numerical weather prediction using Lorenz system benchmarks. \textit{Scientific Reports}, 16, 31979. \url{https://doi.org/10.1038/s41598-026-60719-z}


\bibitem[K\"unsch(1989)]{kunsch1989}
K\"unsch, H. R. (1989). The jackknife and the bootstrap for general stationary observations. \textit{The Annals of Statistics}, 17(3), 1217--1241. \url{https://doi.org/10.1214/aos/1176347265}

\bibitem[Landman et al.(2023)]{landman2023}
Landman, J., Thabet, S., Dalyac, C., Mhiri, H., \& Kashefi, E. (2023). Classically approximating variational quantum machine learning with random Fourier features. In \textit{International Conference on Learning Representations (ICLR 2023)}. \url{https://arxiv.org/abs/2210.13200}

\bibitem[Marchisio et al.(2026)]{marchisio2026}
Marchisio, A., Ebrahim, A., Innan, N., Kashif, M., \& Shafique, M. (2026). QLIF-CAST: Quantum Leaky-Integrate-and-Fire for time-series weather forecasting. \textit{arXiv preprint arXiv:2605.18333}. \url{https://arxiv.org/abs/2605.18333}


\bibitem[Mari et al.(2020)]{mari2020}
Mari, A., Bromley, T. R., Izaac, J., Schuld, M., \& Killoran, N. (2020). Transfer learning in hybrid classical-quantum neural networks. \textit{Quantum}, 4, 340. \url{https://doi.org/10.22331/q-2020-10-28-340}

\bibitem[McClean et al.(2018)]{mcclean2018}
McClean, J. R., Boixo, S., Smelyanskiy, V. N., Babbush, R., \& Neven, H. (2018). Barren plateaus in quantum neural network training landscapes. \textit{Nature Communications}, 9(1), 4812. \url{https://doi.org/10.1038/s41467-018-07090-4}

\bibitem[Miller(2024)]{miller2024}
Miller, E. (2024). Adding error bars to evals: A statistical approach to language model evaluations. arXiv:2411.00640. \url{https://doi.org/10.48550/arXiv.2411.00640}

\bibitem[Nadeau \& Bengio(2003)]{nadeau2003}
Nadeau, C., \& Bengio, Y. (2003). Inference for the generalization error. \textit{Machine Learning}, 52(3), 239--281. \url{https://doi.org/10.1023/A:1024068626366}

\bibitem[Nieto et al.(2005)]{nieto2005}
Nieto, R., Gimeno, L., de la Torre, L., Ribera, P., Gallego, D., Garc\'ia-Herrera, R., Garc\'ia, J. A., Nu\~nez, M., Reda\~no, A., \& Lorente, J. (2005). Climatological features of cutoff low systems in the Northern Hemisphere. \textit{Journal of Climate}, 18(16), 3085--3103. \url{https://doi.org/10.1175/JCLI3386.1}

\bibitem[Obuchowski(1997)]{obuchowski1997}
Obuchowski, N. A. (1997). Nonparametric analysis of clustered ROC curve data. \textit{Biometrics}, 53(2), 567--578. \url{https://doi.org/10.2307/2533958}

\bibitem[Pastori et al.(2026)]{pastori2026}
Pastori, L., Grundner, A., Eyring, V., \& Schwabe, M. (2026). Quantum neural networks for cloud cover parameterizations in climate models. \textit{Machine Learning: Earth}, 2, 015008. \url{https://doi.org/10.1088/3049-4753/ae4981}

\bibitem[Peterson et al.(1998)]{peterson1998}
Peterson, T. C., Easterling, D. R., Karl, T. R., Groisman, P., Nicholls, N., Plummer, N., Torok, S., Auer, I., Boehm, R., Gullett, D., Vincent, L., Heino, R., Tuomenvirta, H., Mestre, O., Szentimrey, T., Salinger, J., F\o rland, E. J., Hanssen-Bauer, I., Alexandersson, H., Jones, P., \& Parker, D. (1998). Homogeneity adjustments of in situ atmospheric climate data: A review. \textit{International Journal of Climatology}, 18(13), 1493--1517. \url{https://doi.org/10.1002/(SICI)1097-0088(19981115)18:13<1493::AID-JOC329>3.0.CO;2-T}

\bibitem[Preskill(2018)]{preskill2018}
Preskill, J. (2018). Quantum computing in the NISQ era and beyond. \textit{Quantum}, 2, 79. \url{https://doi.org/10.22331/q-2018-08-06-79}

\bibitem[Romero et al.(1999)]{romero1999}
Romero, R., Ramis, C., Sumner, G. N., \& Alonso, S. (1999). Daily precipitation patterns in the Spanish Mediterranean area: an objective classification. \textit{International Journal of Climatology}, 19(1), 95--112. \url{https://doi.org/10.1002/(SICI)1097-0088(199901)19:1<95::AID-JOC344>3.0.CO;2-9}

\bibitem[Schuld \& Killoran(2022)]{schuld2022}
Schuld, M., \& Killoran, N. (2022). Is quantum advantage the right goal for quantum machine learning? \textit{PRX Quantum}, 3, 030101. \url{https://doi.org/10.1103/PRXQuantum.3.030101}

\bibitem[Schuld et al.(2021)]{schuld2021fourier}
Schuld, M., Sweke, R., \& Meyer, J. J. (2021). Effect of data encoding on the expressive power of variational quantum machine learning models. \textit{Physical Review A}, 103(3), 032430. \url{https://doi.org/10.1103/PhysRevA.103.032430}

\bibitem[Schreiber et al.(2023)]{schreiber2023}
Schreiber, F. J., Eisert, J., \& Meyer, J. J. (2023). Classical surrogates for quantum learning models. \textit{Physical Review Letters}, 131(10), 100803. \url{https://doi.org/10.1103/PhysRevLett.131.100803}

\bibitem[Schwabe et al.(2025a)]{schwabe2025op}
Schwabe, M., Pastori, L., de Vega, I., Gentine, P., Iapichino, L., Lahtinen, V., Leib, M., Lorenz, J. M., \& Eyring, V. (2025a). Opportunities and challenges of quantum computing for climate modelling. \textit{Environmental Data Science}, 4, e35. \url{https://doi.org/10.1017/eds.2025.10010}

\bibitem[Schwabe et al.(2025b)]{schwabe2025}
Schwabe, M., Pastori, L., Sarandrea, V., \& Eyring, V. (2025b). Quantum machine learning for climate modelling. In \textit{2025 IEEE International Conference on Quantum Artificial Intelligence (QAI)} (pp. 73--78). \url{https://doi.org/10.1109/QAI63978.2025.00019}


\bibitem[Sweke et al.(2025)]{sweke2025}
Sweke, R., Recio-Armengol, E., Jerbi, S., Gil-Fuster, E., Fuller, B., Eisert, J., \& Meyer, J. J. (2025). Potential and limitations of random Fourier features for dequantizing quantum machine learning. \textit{Quantum}, 9, 1640. \url{https://doi.org/10.22331/q-2025-02-20-1640}

\bibitem[Thanasilp et al.(2024)]{thanasilp2024}
Thanasilp, S., Wang, S., Cerezo, M., \& Holmes, Z. (2024). Exponential concentration in quantum kernel methods. \textit{Nature Communications}, 15, 5200. \url{https://doi.org/10.1038/s41467-024-49287-w}

\bibitem[Varoquaux(2018)]{varoquaux2018}
Varoquaux, G. (2018). Cross-validation failure: Small sample sizes lead to large error bars. \textit{NeuroImage}, 180, 68--77. \url{https://doi.org/10.1016/j.neuroimage.2017.06.061}

\bibitem[Wang et al.(2021)]{wang2021noise}
Wang, S., Fontana, E., Cerezo, M., Sharma, K., Sone, A., Cincio, L., \& Coles, P. J. (2021). Noise-induced barren plateaus in variational quantum algorithms. \textit{Nature Communications}, 12, 6961. \url{https://doi.org/10.1038/s41467-021-27045-6}

\bibitem[Webb(2023)]{webb2023}
Webb, M. D. (2023). Reworking wild bootstrap-based inference for clustered errors. \textit{Canadian Journal of Economics}, 56(3), 839--858. \url{https://doi.org/10.1111/caje.12661}

\end{thebibliography}

\appendix

% ============================================================
\section{Computations}
\label{app:noise_budget}
% ============================================================

\textbf{Estimators.} For two models scored on the same 529 rows, the DeLong estimator gives the covariance of the two AUCs from their structural components \citep{delong1988}; it treats station-months as independent. The block bootstrap resamples the 19 calendar years (2008--2026) with replacement, keeping all stations of a year together, recomputes both AUCs on each resample ($B=2{,}000$) and takes percentile intervals \citep{field2007}; the moving-block variant with blocks of several consecutive years \citep{kunsch1989} is a sensitivity check. The two-way bootstrap draws a calendar-year resample and, independently for each model, a resample of its 15 seeds, and recomputes the difference of seed-averaged AUCs. The cluster-robust estimator uses the influence values $\psi_i=(\bar v^{A}_{10,i}-\bar v^{B}_{10,i}-\hat\Delta)/n_1$ for positives and $(\bar v^{A}_{01,j}-\bar v^{B}_{01,j}-\hat\Delta)/n_0$ for negatives, where $\bar v$ are the DeLong components averaged over seeds; $\sum_i\psi_i^2$ equals the DeLong variance of $\hat\Delta$ up to factors $(n-1)/n$, and the year sums of $\psi$ give the cluster-robust one. Writing the year sum as $S_u=\sum_s S_{u,s}$ over stations $s$, $\sum_u S_u^2=\sum_{u,s}S_{u,s}^2+\sum_u\sum_{s\ne t}S_{u,s}S_{u,t}$ separates the within-station and cross-station terms; the permutation null permutes the year labels independently within each station (2{,}000 permutations).

\textbf{Seed-mean statistic.} The AUC of one model is the mean of the DeLong structural components over positives and over negatives, so the mean of the 15 AUCs equals the AUC statistic built from components averaged over seeds, and the paired DeLong variance of $\hat\Delta$ follows from the averaged components. This treats the 15 seeds as fixed and so conditions on the fitted models; the two-way bootstrap relaxes this by resampling seeds.

\textbf{Simulation mechanism.} For each replication, years $u=1,\dots,19$ receive a prevalence effect $\tau_\pi\,\eta_u$ on the logit of $\pi=0.41$ and a number of rows of 528 split as equally as possible (or, in the unequal scenario, a multinomial split with Dirichlet$(3)$ shares). Scorers $k\in\{A,B\}$ are $s_{k}=\mu_k y+e_k+\tau_s\,w_u+\tau_m\,v_{k,u}$, with $e_k=\sqrt\rho\,z+\sqrt{1-\rho}\,\varepsilon_k$ (so that the scorer noise has correlation $\rho$), $w_u$ a year shift shared by both scorers and $v_{k,u}$ a year shift specific to each scorer; all $\eta_u,w_u,v_{k,u},z,\varepsilon_k$ are standard normal and $\mu_k=\sqrt2\,\Phi^{-1}(\mathrm{AUC}_k)$ without year effects. The design AUCs without year effects are 0.65 for scorer B and $0.65+\Delta$ for scorer A; the scorer noise has correlation $\rho=0.9$, and the scorer-specific year shifts $v_{k,u}$ are independent between scorers. Scenario B has $\tau_\pi=\tau_s=0.3$ and $\tau_m=0$; C adds $\tau_m=0.3$; D is C with Dirichlet year sizes; E has $\tau_\pi=\tau_s=\tau_m=0.6$. The simulation target is the pooled population difference under the mechanism, with years drawn independently and weighted equally, that is, the difference for new years; in scenario D the shares of the 528 rows among years follow a Dirichlet(3) draw that is independent of the year effects, the sample statistic pools observations (so it weights years by their sizes), and the reference is the pooled population AUC with equally likely years of equal size. Scenario D changes the distribution of cluster sizes. Because the sizes are independent of the year effects, the pooled statistic is expected to be centered on the equal-weight reference; this was checked empirically and not proved: in D the mean of the estimates differs from the reference by $+0.0002$, $-0.0008$ and $-0.0004$ for design differences of 0, 0.015 and 0.03 (Monte Carlo standard error of the mean about 0.0008). The population AUC of each scorer is computed as the mean of four independent draws of 400{,}000 rows from the same mechanism with years drawn independently and equally weighted, and the Monte Carlo standard error of the population difference is reported. When the design difference is 0 the scorers are exchangeable and the population difference is set to 0. DeLong tests use the paired estimator of the seed-free case; the block bootstrap resamples years ($B=300$ per replication).

\end{document}
